% \iffalse
%  Local Variables:
%  mode: doctex
%  TeX-master: t
%  End:
% \fi
%
% \iffalse meta-comment
%
% Copyright (C) 2005-2013 by Ruini Xue <xueruini@gmail.com>
% Copyright (C) 2013-2018 derived by Chaolu <wuchaolu@gmail.com>
% This file is part of the UCASThesis package project, which is
% originally derived from ThuThesis under LPPL.
%
% This file may be distributed and/or modified under the
% conditions of the LaTeX Project Public License, either version 1.3a
% of this license or (at your option) any later version.
% The latest version of this license is in:
%
% http://www.latex-project.org/lppl.txt
%
% and version 1.3a or later is part of all distributions of LaTeX
% version 2004/10/01 or later.
%
% $Id$
%
% \fi
%
% \CheckSum{0}
% \CharacterTable
%  {Upper-case    \A\B\C\D\E\F\G\H\I\J\K\L\M\N\O\P\Q\R\S\T\U\V\W\X\Y\Z
%   Lower-case    \a\b\c\d\e\f\g\h\i\j\k\l\m\n\o\p\q\r\s\t\u\v\w\x\y\z
%   Digits        \0\1\2\3\4\5\6\7\8\9
%   Exclamation   \!     Double quote  \"     Hash (number) \#
%   Dollar        \$     Percent       \%     Ampersand     \&
%   Acute accent  \'     Left paren    \(     Right paren   \)
%   Asterisk      \*     Plus          \+     Comma         \,
%   Minus         \-     Point         \.     Solidus       \/
%   Colon         \:     Semicolon     \;     Less than     \<
%   Equals        \=     Greater than  \>     Question mark \?
%   Commercial at \@     Left bracket  \[     Backslash     \\
%   Right bracket \]     Circumflex    \^     Underscore    \_
%   Grave accent  \`     Left brace    \{     Vertical bar  \|
%   Right brace   \}     Tilde         \~}
%
% \iffalse
%<*driver>
\ProvidesFile{ucasthesis.dtx}[2014/10/01 1.0 dev University of Chinese Academy of Sciences Thesis Template]
\documentclass[10pt]{ltxdoc}
\usepackage{dtx-style}
\EnableCrossrefs
\CodelineIndex
\RecordChanges
%\OnlyDescription
\begin{document}
  \DocInput{\jobname.dtx}
\end{document}
%</driver>
% \fi
%
% \GetFileInfo{\jobname.dtx}
% \MakeShortVerb{\|}
%
% \def\ucasthesis{\textsc{ucas}\-\textsc{Thesis}}
% \def\pkg#1{\texttt{#1}}
%
%
% \DoNotIndex{\begin,\end,\begingroup,\endgroup}
% \DoNotIndex{\ifx,\ifdim,\ifnum,\ifcase,\else,\or,\fi}
% \DoNotIndex{\let,\def,\xdef,\newcommand,\renewcommand}
% \DoNotIndex{\expandafter,\csname,\endcsname,\relax,\protect}
% \DoNotIndex{\Huge,\huge,\LARGE,\Large,\large,\normalsize}
% \DoNotIndex{\small,\footnotesize,\scriptsize,\tiny}
% \DoNotIndex{\normalfont,\bfseries,\slshape,\interlinepenalty}
% \DoNotIndex{\hfil,\par,\hskip,\vskip,\vspace,\quad}
% \DoNotIndex{\centering,\raggedright}
% \DoNotIndex{\c@secnumdepth,\@startsection,\@setfontsize}
% \DoNotIndex{\ ,\@plus,\@minus,\p@,\z@,\@m,\@M,\@ne,\m@ne}
% \DoNotIndex{\@@par,\DeclareOperation,\RequirePackage,\LoadClass}
% \DoNotIndex{\AtBeginDocument,\AtEndDocument}
%
% \IndexPrologue{\section*{索引}%
%    \addcontentsline{toc}{section}{索~~~~引}}
% \GlossaryPrologue{\section*{修改记录}%
%    \addcontentsline{toc}{section}{修改记录}}
%
% \renewcommand{\abstractname}{摘~~要}
% \renewcommand{\contentsname}{目~~录}
%
% \makeatletter
% \let\@fnsymbol\@arabic
% \makeatother
%
% \title{\ucasthesis：中国科学院大学学位论文模板\thanks{University of Chinese Academy of Sciences \LaTeX{} Thesis Template.}}
% \author{{\fangsong 作者：朝鲁\thanks{wuchaolu@gmail.com}}\\[5pt]{\fangsong 中国科学院计算技术研究所}\\[5pt] \texttt{wuchaolu@gmail.com}\\{\fangsong 原作者：薛瑞尼\thanks{LittleLeo@newsmth}}\\[5pt]{\fangsong 清华大学计算机系高性能所}\\[5pt] \texttt{xueruini@gmail.com}}
% \date{v\fileversion\ (\filedate)}
% \maketitle\thispagestyle{empty}
% \changes{v1.0 }{2014/09/20}{Rebase from ThuThesis v4.8。为防止后续开发工作混淆，在UCASThesis中清空了过往的changes
% ，如有需要查阅先前changelog的开发者，请参照ThuThesis的github log。本版本重新适应了中国科学院大学学位论文需求，并重新
% 定义版本号为v1.0。论文模板的一些针对清华大学学位论文文档格式的细致说明问题，由于我本人时间精力有限，目前并未进行较大
% 改动，请各位见谅。}
% \changes{v1.1 }{2016/01/05}{向上兼容支持Tex2015, CTEX2.0+，同时开始使用固定字体配置。}
% \changes{v1.2 }{2018/03/05}{替换了BST文件为胡海星的版本，支持了online以及多种文献格式。同时更新了相关版权信息。}
% \changes{v2.0 }{2018/03/17}{更新模板为国科大统一标准。}

%
% \begin{abstract}\noindent
%   此宏包旨在建立一个简单易用的中国科学院大学学位论文模板，包括硕士
%   论文、博士论文。
% \end{abstract}
%
% \vskip2cm
% \def\abstractname{免责声明}
% \begin{abstract}
% \noindent
% \begin{enumerate}
% \item 本模板是以薛瑞尼维护的清华大学学位论文模板（ThuThesis）为基础制作的衍生版，
%  本模板发布遵守 \LaTeX{} Project Public License，使用前请认真阅读协议内容。
% \item 本模板依据《中国科学院大学研究生学位论文撰写规定》、《国科大版论文模板》，
%  并参考中科院计算所提供的Word版学位论文示例的具体格式制作而成，
%  旨在供中国科学院大学毕业生撰写学位论文使用。
% \item 本模板仅为作者个人对官方文件的参考实现，不保证中国科学院大学或各个研究所负责格式
%  审查的老师不提意见。任何由于使用本模板而引起的论文格式审查问题均与本模板作者无关。
% \item 任何个人或组织以本模板为基础进行修改、扩展而生成的新的专用模板，请严格遵
%  守 \LaTeX{} Project Public License 协议。由于违反协议而引起的任何纠纷争端均与
%  本模板作者无关。
% \end{enumerate}
% \end{abstract}
%
%
% \clearpage
% \begin{multicols}{2}[
%   \section*{\contentsname}
%   \setlength{\columnseprule}{.4pt}
%   \setlength{\columnsep}{18pt}]
%   \tableofcontents
% \end{multicols}
%
% \clearpage
% \pagenumbering{arabic}
% \pagestyle{headings}
% \section{模板介绍}
% \ucasthesis\ (\textbf{UCAS} \textbf{Thesis}) 是为了帮助中国科学院大学毕业
% 生撰写学位论文而编写的 \LaTeX{} 论文模板。
%
% 本文档将尽量完整的介绍模板的使用方法，如有不清楚之处可以参考示例文档或者给邮件
% 列表（见后）写信，欢迎感兴趣的同学出力完善此使用手册。由于个人水平有限，虽然现
% 在的这个版本基本上满足了学校的要求，但难免还存在不足之处，欢迎大家积极反馈。
%
% {\color{blue}\fangsong 模板的作用在于减轻论文写作过程中格式调整的时间，其前提就是遵
%   守模板的用法，否则即使使用了 \ucasthesis{} 也难以保证输出的论文符合学校规范。}
%
%
% \section{安装}
% \label{sec:installation}
%
% \subsection{ThuThesis下载}
% ThuThesis 相关链接：
% \begin{itemize}
% \item 主页：
% \href{https://github.com/xueruini/thuthesis}{Github}\footnote{已经从
% \url{http://thuthesis.sourceforge.net}迁移至此。}
% \item 下载：\href{http://code.google.com/p/thuthesis/}{Google Code}
% \item 同时本模板也提交至
% \href{http://www.ctan.org/macros/latex/contrib/thuthesis}{CTAN}
% \end{itemize}
% 除此之外，不再维护任何镜像。
%
% ThuThesis 的开发版本同样可以在 Github 上获得：
% \begin{shell}
% $ git clone git://github.com/xueruini/thuthesis.git
% \end{shell}
% 
% \subsection{UCASThesis下载}
% UCASThesis 相关链接：
% \href{https://github.com/xiaoyao9933/ucasthesis}{Github}
% UCASThesis 的开发版本同样可以在 Github 上获得：
% \begin{shell}
% $ git clone git://github.com/xiaoyao9933/ucasthesis.git
% \end{shell}
% \subsection{模板的组成部分}
% 下表列出了 \ucasthesis{} 的主要文件及其功能介绍：
%
% \begin{center}
%   \begin{longtable}{l|p{10cm}}
% \hline
% {\heiti 文件（夹）} & {\heiti 功能描述}\\\hline\hline
% \endfirsthead
% \hline
% {\heiti 文件（夹）} & {\heiti 功能描述}\\\hline\hline
% \endhead
% \endfoot
% \endlastfoot
% ucasthesis.ins & 模板驱动文件 \\
% ucasthesis.dtx & 模板文档代码的混合文件\\
% ucasthesis.cls & 模板类文件\\
% ucasthesis.cfg & 模板配置文件\\
% ucasbib.bst & 参考文献样式文件\\\hline
% main.tex & 示例文档主文件\\
% spine.tex & 书脊示例文档\\
% ref/ & 示例文档参考文献目录\\
% data/ & 示例文档章节具体内容\\
% figures/ & 示例文档图片路径\\
% ucastils.sty & 为示例文档加载其它宏包\\\hline
% Makefile & self-explanation \\
% Readme & self-explanation\\
% \textbf{ucasthesis.pdf} & 用户手册（本文档）\\\hline
%   \end{longtable}
% \end{center}
%
% 需要说明几点：
% \begin{itemize}
% \item \emph{ucasthesis.cls} 和 \emph{ucasthesis.cfg} 可以
%   由 \emph{ucasthesis.ins} 和 \emph{ucasthesis.dtx} 生成，但为了降低新
%   手用户的使用难度，故将 cls和 cfg 一起发布。
% \item 使用前认真阅读文档：\emph{ucasthesis.pdf}.
% \end{itemize}
% 
% \subsection{准备工作}
% \label{sec:prepare}
% 本模板需要基于texlive 2014以上版本编译通过。\\
% 本模板用到以下宏包：
%
% \begin{center}
% \begin{minipage}{1.0\linewidth}\centering
% \begin{tabular}{*{6}{l}}\hline
%   ifxetex & xunicode & CJK\footnote{版本要求：$\geq$ v4.8.1} & xeCJK & \pkg{CJKpunct} & \pkg{ctex} \\
%   array & booktabs & longtable  &  amsmath & amssymb & ntheorem \\
%   indentfirst & paralist & txfonts & natbib & hyperref \\
%   graphicx & \pkg{subcaption} &
%   \pkg{caption}\footnote{版本要求：$\geq$2006/03/21 v3.0j} &
%   \pkg{ucasbib.bst} & &\\\hline
% \end{tabular}
% \end{minipage}
% \end{center}
%
% 这些包在常见的 \TeX{} 系统中都有，如果没有请到 \url{www.ctan.org} 下载。推
% 荐 \TeX\ Live。
%
%
% \subsection{开始安装}
% \label{sec:install}
%
% \subsubsection{生成模板}
% \label{sec:generate-cls}
% {\heiti 说明：默认的发行包中已经包含了所有文件，可以直接使用。如果对如何由 dtx 生
%   成模板文件以及模板文档不感兴趣，请跳过本小节。}
%
% 模板解压缩后生成文件夹 ucasthesis-VERSION\footnote{VERSION 为版本号。}，其中包括：
% 模板源文件（ucasthesis.ins 和 ucasthesis.dtx），参考文献样式 ucasbib.bst，示例文档
% （main.tex，spine.tex，ucastils.sty\footnote{我把可能用到但不一定用到的包以及一
%   些命令定义都放在这里面，以免 ucasthesis.cls 过分臃
%   肿。}，data/ 和 figures/ 和 ref/）。在使用之前需要先生成模板文件和配置文件
% （具体命令细节请参考 |Readme| 和 |Makefile|）：
%
% \begin{shell}
% $ cd ucasthesis-VERSION
% # 生成 ucasthesis.cls 和 ucasthesis.cfg
% $ latex ucasthesis.ins
%
% # 下面的命令用来生成用户手册，可以不执行
% $ latex ucasthesis.dtx
% $ makeindex -s gind.ist -o ucasthesis.ind ucasthesis.idx
% $ makeindex -s gglo.ist -o ucasthesis.gls ucasthesis.glo
% $ latex ucasthesis.dtx
% $ latex ucasthesis.dtx  % 生成说明文档 ucasthesis.dvi
% \end{shell}
%
%
% \subsubsection{dvi$\rightarrow$ps$\rightarrow$pdf}
% \label{sec:dvipspdf}
% 很多用户对 \LaTeX{} 命令执行的次数不太清楚，一个基本的原则是多次运行 \LaTeX{}
% 命令直至不再出现警告。下面给出生成示例文档的详细过程（\# 开头的行为注释），首先
% 来看经典的 \texttt{dvi$\rightarrow$ps$\rightarrow$pdf} 方式：
% \begin{shell}
% # 1. 发现里面的引用关系，文件后缀 .tex 可以省略
% $ latex main
%
% # 2. 编译参考文件源文件，生成 bbl 文件
% $ bibtex main
%
% # 3. 下面解决引用
% $ latex main
% # 如果是 GBK 编码，此处运行：
% # $ gbk2uni main  # 防止书签乱码
% $ latex main   # 此时生成完整的 dvi 文件
%
% # 4. 生成 ps
% $ dvips main.dvi
%
% # 5. 生成 pdf
% $ ps2pdf main.ps
% \end{shell}
%
% 模板已经把纸型信息写入目标文件，这样执行 \texttt{dvips} 时就可以避免由于遗忘
%  \texttt{-ta4} 参数而导致输出不合格的文件（因为 \texttt{dvips} 默认使用
%  letter 纸型）。
%
% \subsubsection{dvipdfm(x)}
% \label{sec:dvipdfmx}
% 如果使用 dvipdfm(x)，那么在生成完整的 dvi 文件之后（参见上面的例子），可以直接得到 pdf：
% \begin{shell}%
% $ dvipdfm  main.dvi
% # 或者
% $ dvipdfmx  main.dvi
% \end{shell}
%
% \subsubsection{pdflatex}
% \label{sec:pdflatex}
% 如果使用 PDF\LaTeX，按照第~\ref{sec:dvipspdf} 节的顺序执行到第 3 步即可，不再经
% 过中间转换。
%
% 需要注意的是 PDF\LaTeX\ 不能处理常见的 EPS 图形，需要先用 epstopdf 将其转化
% 成 PDF。不过 PDF\LaTeX\ 增加了对 png，jpg 等标量图形的支持，比较方便。
%
% \subsubsection{xelatex}
% \label{sec:xelatex}
% XeTeX 最大的优势就是不再需要繁琐的字体配置。\ucasthesis{} 通过 \pkg{xeCJK} 来控
% 制中文字体和标点压缩。模板里默认用的是中易的四款免费字体（宋，黑，楷，仿宋），
% 用户可以根据自己的实际情况方便的替换。
%
% Xe\LaTeX\ 的使用步骤同 PDF\LaTeX。
%
%
% \subsubsection{自动化过程}
% \label{sec:automation}
% 上面的例子只是给出一般情况下的使用方法，可以发现虽然命令很简单，但是每次都输入
% 的话还是非常罗嗦的，所以 \ucasthesis{} 还提供了一些自动处理的文件。
%
% 我们提供了一个简单的 \texttt{Makefile}：
% \begin{shell}
% $ make clean
% $ make cls       # 生成 ucasthesis.cls 和 ucasthesis.cfg
% $ make doc       # 生成说明文档 ucasthesis.pdf
% $ make thesis    # 生成示例文档 main.pdf
% $ make spine     # 生成书脊 spine.pdf
% \end{shell}
%
% \texttt{Makefile} 默认采用 Xe\LaTeX\ 编译，可以根据自己的
% 需要修改 \texttt{config.mk} 中的参数设置。
%
%
% \subsection{升级}
% \label{sec:updgrade}
% \ucasthesis{} 升级非常简单，下载最新的版本，
% 将 ucasthesis.ins，ucasthesis.dtx 和ucasbib.bst 拷贝至工作目录覆盖相应的文件，然后
% 运行：
% \begin{shell}
% $ latex ucasthesis.ins
% \end{shell}
%
% 生成新的类文件和配置文件即可。当然也可以直接拷贝 ucasthesis.cls, ucasthesis.cfg
% 和 ucasbib.bst，免去上面命令的执行。只要明白它的工作原理，这个不难操作。
%
%
% \section{使用说明}
% \label{sec:usage}
% 本手册假定用户已经能处理一般的 \LaTeX{} 文档，并对 \BibTeX{} 有一定了解。如果你
% 从来没有接触过 \TeX 和 \LaTeX，建议先学习相关的基础知识。磨刀不误砍柴工！
%
% \subsection{关于提问}
% \label{sec:howtoask}
% \begin{itemize}\addtolength{\itemsep}{-5pt}
% \item \url{http://groups.google.com/group/thuthesis}
% 或直接给\href{mailto:thuthesis@googlegroups.com}{邮件列表}写信。
% \item Google Groups mirror: \url{http://thuthesis.1048723.n5.nabble.com/}
% \item \href{http://www.newsmth.net/bbsdoc.php?board=TeX}{\TeX@newsmth}
% \end{itemize}
%
% \subsection{\ucasthesis{} 使用向导}
% \label{sec:userguide}
% 推荐新用户先看网上的《\ucasthesis{} 使用向导》幻灯片\footnote{有点老了，不过还是
%   很有帮助的。}，那份讲稿比这份文档简练易懂。
%
% \subsection{\ucasthesis{} 示例文件}
% \label{sec:userguide1}
% 模板核心文件只有三个：ucasthesis.cls，ucasthesis.cfg 和 ucasbib.bst，但是如果没有
% 示例文档用户会发现很难下手。所以推荐新用户从模板自带的示例文档入手，里面包括了
% 论文写作用到的所有命令及其使用方法，只需要用自己的内容进行相应替换就可以。对于
% 不清楚的命令可以查阅本手册。下面的例子描述了模板中章节的组织形式，来自于示例文
% 档，具体内容可以参考模板附带的 main.tex 和 data/。
%
% \begin{example}
% 
% \documentclass[master,adobefonts]{ucasthesis}
% %\documentclass[doctor]{ucasthesis}
% %\documentclass[%
% %  master|doctor, % 必选选项
% %  secret, % 可选选项
% %  openany|openright, % 可选选项
% %  arialtoc,arialtitle % 可选选项
% %  ]{ucasthesis}
%
% % 所有其它可能用到的包都统一放到这里了，可以根据自己的实际添加或者删除。
% \usepackage{ucastils}
%
% % 可以在这里修改配置文件中的定义，导言区可以使用中文。
% % \def\myname{薛瑞尼}
%
% \begin{document}
%
% % 指定图片的搜索目录
% \graphicspath{{figures/}}
%
%
% %%% 封面部分
% \frontmatter
% 
%%% Local Variables:
%%% mode: latex
%%% TeX-master: t
%%% End:
\secretcontent{}  % 将这里的括号内文字修改为空白即可保持下划线，同时没有任何文字显示。

\ctitle{中国科学院大学学位论文 \LaTeX 模板\\使用示例文档}
% 根据自己的情况选，不用这样复杂
\makeatletter

\makeatother

\cdegree{工学博士}  %
\cdepartment[计算所]{中国科学院计算技术研究所}
\cmajor{计算机系统结构}  % 二级学科专业名称，如遗传学等
\cauthor{朝\hspace{1em}鲁} 
\csupervisor{徐志伟\hspace{1em}研究员}
\csupervisorplace{中国科学院计算技术研究所}
% 如果没有副指导老师或者联合指导老师，把下面两行相应的删除即可。


% 日期自动生成，如果你要自己写就改这个cdate
%\cdate{2018年6月} % 年月用阿拉伯数字填写，Times New Roman体四号加粗，夏季毕业填写6月，冬季填写12月

\etitle{An Introduction to \LaTeX{} Thesis Template of\\University of Chinese Academy of Sciences} 

\edegree{Doctor of Philosophy}

\eauthor{Chao Lu}
\edepartment{Institute of Computing Technology\\Chinese Academy of Sciences}
\emajor{XXXXXXXXXXXX}   %专业名称，具体填一级专业和二级专业，请咨询你们的老师
\esupervisor{Xu Zhiwei}

% 这个日期也会自动生成，你要改么？
% \edate{December 2016}

% 定义中英文摘要和关键字
\begin{cabstract}
  论文的摘要是对论文研究内容和成果的高度概括。摘要应对论文所研究的问题及其研究目
  的进行描述，对研究方法和过程进行简单介绍，对研究成果和所得结论进行概括。摘要应
  具有独立性和自明性，其内容应包含与论文全文同等量的主要信息。使读者即使不阅读全
  文，通过摘要就能了解论文的总体内容和主要成果。

  论文摘要的书写应力求精确、简明。切忌写成对论文书写内容进行提要的形式，尤其要避
  免“第 1 章……；第 2 章……；……”这种或类似的陈述方式。

  本文介绍中国科学院大学论文模板 \ucasthesis{} 的使用方法。本模板符合学校的硕士、
  博士论文格式要求。

  本文的创新点主要有：
  \begin{itemize}
    \item 用例子来解释模板的使用方法；
    \item 用废话来填充无关紧要的部分；
    \item 一边学习摸索一边编写新代码。
  \end{itemize}

  关键词是为了文献标引工作、用以表示全文主要内容信息的单词或术语。关键词不超过 5
  个，每个关键词中间用分号分隔。（模板作者注：关键词分隔符不用考虑，模板会自动处
  理。英文关键词同理。）
\end{cabstract}

\ckeywords{\TeX, \LaTeX, CJK, 模板, 论文}

\begin{eabstract} 
   An abstract of a dissertation is a summary and extraction of research work
   and contributions. Included in an abstract should be description of research
   topic and research objective, brief introduction to methodology and research
   process, and summarization of conclusion and contributions of the
   research. An abstract should be characterized by independence and clarity and
   carry identical information with the dissertation. It should be such that the
   general idea and major contributions of the dissertation are conveyed without
   reading the dissertation. 

   An abstract should be concise and to the point. It is a misunderstanding to
   make an abstract an outline of the dissertation and words ``the first
   chapter'', ``the second chapter'' and the like should be avoided in the
   abstract.

   Key words are terms used in a dissertation for indexing, reflecting core
   information of the dissertation. An abstract may contain a maximum of 5 key
   words, with semi-colons used in between to separate one another.
\end{eabstract}

\ekeywords{\TeX, \LaTeX, CJK, template, thesis}

% \makecover
%
% % 目录
% \tableofcontents
%
% % 符号对照表
% \input{data/denotation}
%
%
% %%% 正文部分
% \mainmatter
% 
%%% Local Variables:
%%% mode: latex
%%% TeX-master: t
%%% End:

\chapter{带 English 的标题}
\label{cha:intro}

这是 \ucasthesis{} 的示例文档，基本上覆盖了模板中所有格式的设置。建议大家在使用模
板之前，除了阅读《\ucasthesis{}用户手册》，这个示例文档也最好能看一看。

小老鼠偷吃热凉粉；短长虫环绕矮高粱。\footnote{韩愈（768-824），字退之，河南河阳（
  今河南孟县）人，自称郡望昌黎，世称韩昌黎。幼孤贫刻苦好学，德宗贞元八年进士。曾
  任监察御史，因上疏请免关中赋役，贬为阳山县令。后随宰相裴度平定淮西迁刑部侍郎，
  又因上表谏迎佛骨，贬潮州刺史。做过吏部侍郎，死谥文公，故世称韩吏部、韩文公。是
  唐代古文运动领袖，与柳宗元合称韩柳。诗力求险怪新奇，雄浑重气势。}
  
\section{封面相关}
封面的例子请参看 cover.tex。主要符号表参看 denation.tex，附录和个人简历分别参看 appendix01.tex
和 resume.tex。里面的命令都非常简单，一看即会。\footnote{你说还是看不懂？怎么会呢？}

\section{字体命令}
\label{sec:first}

苏轼（1037-1101），北宋文学家、书画家。字子瞻，号东坡居士，眉州眉山（今属四川）人
。苏洵子。嘉佑进士。神宗时曾任祠部员外郎，因反对王安石新法而求外职，任杭州通判，
知密州、徐州、湖州。后以作诗“谤讪朝廷”罪贬黄州。哲宗时任翰林学士，曾出知杭州、
颖州等，官至礼部尚书。后又贬谪惠州、儋州。北还后第二年病死常州。南宋时追谥文忠。
与父洵弟辙，合称“三苏”。在政治上属于旧党，但也有改革弊政的要求。其文汪洋恣肆，
明白畅达，为“唐宋八大家”之一。  其诗清新豪健，善用夸张比喻，在艺术表现方面独具
风格。少数诗篇也能反映民间疾苦，指责统治者的奢侈骄纵。词开豪放一派，对后代很有影
响。《念奴娇·赤壁怀古》、《水调歌头·丙辰中秋》传诵甚广。

{\kaishu 坡仙擅长行书、楷书，取法李邕、徐浩、颜真卿、杨凝式，而能自创新意。用笔丰腴
  跌宕，有天真烂漫之趣。与蔡襄、黄庭坚、米芾并称“宋四家”。能画竹，学文同，也喜
  作枯木怪石。论画主张“神似”，认为“论画以形似，见与儿童邻”；高度评价“诗中有
  画，画中有诗”的艺术造诣。诗文有《东坡七集》等。存世书迹有《答谢民师论文帖》、
  《祭黄几道文》、《前赤壁赋》、《黄州寒食诗帖》等。  画迹有《枯木怪石图》、《
  竹石图》等。}

{\fangsong 易与天地准，故能弥纶天地之道。仰以观於天文，俯以察於地理，是故知幽明之故。原
  始反终，故知死生之说。精气为物，游魂为变，是故知鬼神之情状。与天地相似，故不违。
  知周乎万物，而道济天下，故不过。旁行而不流，乐天知命，故不忧。安土敦乎仁，故
  能爱。范围天地之化而不过，曲成万物而不遗，通乎昼夜之道而知，故神无方而易无体。}

% 非本科生一般用不到幼圆与隶书字体。ctex 在 xelatex 编译时用 winfonts/adobefonts
% 选项也只配置了四款中文字体，没有提供幼圆和隶书。需要的同学可以使用 nofonts 选项
% 自行配置中文字体，或者换用 pdflatex 引擎编译。
{\ifcsname youyuan\endcsname\youyuan 有天地，然后万物生焉。盈天地之间者，唯万物，故受之以屯；屯者盈也，屯者物之
  始生也。物生必蒙，故受之以蒙；蒙者蒙也，物之穉也。物穉不可不养也，故受之以需；
  需者饮食之道也。饮食必有讼，故受之以讼。讼必有众起，故受之以师；师者众也。众必
  有所比，故受之以比；比者比也。比必有所畜也，故受之以小畜。物畜然后有礼，故受之
  以履。\fi}

{\heiti 履而泰，然后安，故受之以泰；泰者通也。物不可以终通，故受之以否。物不可以终
  否，故受之以同人。与人同者，物必归焉，故受之以大有。有大者不可以盈，故受之以谦。
  有大而能谦，必豫，故受之以豫。豫必有随，故受之以随。以喜随人者，必有事，故受
  之以蛊；蛊者事也。}

{\ifcsname lishu\endcsname\lishu 有事而后可大，故受之以临；临者大也。物大然后可观，故受之以观。可观而后有所合
  ，故受之以噬嗑；嗑者合也。物不可以苟合而已，故受之以贲；贲者饰也。致饰然后亨
  ，则尽矣，故受之以剥；剥者剥也。物不可以终尽，剥穷上反下，故受之以复。复则不
  妄矣，故受之以无妄。\fi}

{\songti 有无妄然后可畜，故受之以大畜。物畜然后可养，故受之以颐；颐者养也。不养则不
  可动，故受之以大过。物不可以终过，故受之以坎；坎者陷也。陷必有所丽，故受之以
  离；离者丽也。}

\section{表格样本}
\label{chap1:sample:table} 

\subsection{基本表格}
\label{sec:basictable}

模板中关于表格的宏包有三个： \textsf{booktabs}、\textsf{array} 和
\textsf{longtabular}，命令有一个 \verb|\hlinewd|。三线表可以用 \textsf{booktabs}
提供的 \verb|\toprule|、\verb|\midrule| 和 \verb|\bottomrule|。它们与
\textsf{longtable} 能很好的配合使用。如果表格比较简单的话可以直接用命令
\verb|hlinewd{xpt}| 控制。
\begin{table}[htb]
  \centering
  \begin{minipage}[t]{0.8\linewidth} % 如果想在表格中使用脚注，minipage是个不错的办法
  \bicaption[模板文件]{模板文件。如果表格的标题很长，那么在表格索引中就会很不美
    观，所以要像 chapter 那样在前面用中括号写一个简短的标题。这个标题会出现在索
    引中。}[Template File]{Template File}
  \label{tab:template-files}
    \begin{tabular*}{\linewidth}{lp{10cm}}
      \toprule[1.5pt]
      {\heiti 文件名} & {\heiti 描述} \\\midrule[1pt]
      ucasthesis.ins & \LaTeX{} 安装文件，docstrip\footnote{表格中的脚注} \\
      ucasthesis.dtx & 所有的一切都在这里面\footnote{再来一个}。\\
      ucasthesis.cls & 模板类文件。\\
      ucasthesis.cfg & 模板配置文。cls 和 cfg 由前两个文件生成。\\
      ucasbib.bst    & 参考文献 Bibtex 样式文件。\\
      ucastils.sty   & 常用的包和命令写在这里，减轻主文件的负担。\\
      \bottomrule[1.5pt]
    \end{tabular*}
  \end{minipage}
\end{table}

首先来看一个最简单的表格。表 \ref{tab:template-files} 列举了本模板主要文件及其功
能。请大家注意三线表中各条线对应的命令。这个例子还展示了如何在表格中正确使用脚注。
由于 \LaTeX{} 本身不支持在表格中使用 \verb|\footnote|，所以我们不得不将表格放在
小页中，而且最好将表格的宽度设置为小页的宽度，这样脚注看起来才更美观。

\subsection{复杂表格}
\label{sec:complicatedtable}

我们经常会在表格下方标注数据来源，或者对表格里面的条目进行解释。前面的脚注是一种
不错的方法，如果你不喜欢脚注。那么完全可以在表格后面自己写注释，比如表~\ref{tab:tabexamp1}。
\begin{table}[htbp]
  \centering
  \bicaption{复杂表格示例 1}{Complex Table Demo 1}
  \label{tab:tabexamp1}
  \begin{minipage}[t]{0.8\textwidth} 
    \begin{tabularx}{\linewidth}{|l|X|X|X|X|}
      \hline
 \multirow{2}*{\diagbox[width=5em]{x}{y}}  & \multicolumn{2}{c|}{First Half} & \multicolumn{2}{c|}{Second Half}\\\cline{2-5}
      & 1st Qtr &2nd Qtr&3rd Qtr&4th Qtr \\ \hline
      East$^{*}$ &   20.4&   27.4&   90&     20.4 \\
      West$^{**}$ &   30.6 &   38.6 &   34.6 &  31.6 \\ \hline
    \end{tabularx}\\[2pt]
    \footnotesize 注：数据来源《\ucasthesis{} 使用手册》。\\
    *：东部\\
    **：西部
  \end{minipage}
\end{table}

此外，表~\ref{tab:tabexamp1} 同时还演示了另外两个功能：1）通过 \textsf{tabularx} 的
 \texttt{|X|} 扩展实现表格自动放大；2）通过命令 \verb|\diagbox| 在表头部分
插入反斜线。

为了使我们的例子更接近实际情况，我会在必要的时候插入一些“无关”文字，以免太多图
表同时出现，导致排版效果不太理想。第一个出场的当然是我的最爱：风流潇洒、骏马绝尘、
健笔凌云的{\heiti 李太白}了。

李白，字太白，陇西成纪人。凉武昭王暠九世孙。或曰山东人，或曰蜀人。白少有逸才，志
气宏放，飘然有超世之心。初隐岷山，益州长史苏颋见而异之，曰：“是子天才英特，可比
相如。”天宝初，至长安，往见贺知章。知章见其文，叹曰：“子谪仙人也。”言于明皇，
召见金銮殿，奏颂一篇。帝赐食，亲为调羹，有诏供奉翰林。白犹与酒徒饮于市，帝坐沉香
亭子，意有所感，欲得白为乐章，召入，而白已醉。左右以水颒面，稍解，援笔成文，婉丽
精切。帝爱其才，数宴见。白常侍帝，醉，使高力士脱靴。力士素贵，耻之，摘其诗以激杨
贵妃。帝欲官白，妃辄沮止。白自知不为亲近所容，恳求还山。帝赐金放还。乃浪迹江湖，
终日沉饮。永王璘都督江陵，辟为僚佐。璘谋乱，兵败，白坐长流夜郎，会赦得还。族人阳
冰为当涂令，白往依之。代宗立，以左拾遗召，而白已卒。文宗时，诏以白歌诗、裴旻剑舞、
张旭草书为三绝云。集三十卷。今编诗二十五卷。\hfill ——《全唐诗》诗人小传

浮动体的并排放置一般有两种情况：1）二者没有关系，为两个独立的浮动体；2）二者隶属
于同一个浮动体。对表格来说并排表格既可以像图~\ref{tab:parallel1}、图~\ref{tab:parallel2} 
使用小页环境，也可以如图~\ref{tab:subtable} 使用子表格来做。图的例子参见第~\ref{sec:multifig} 节。
\begin{table}[htbp]
\noindent\begin{minipage}{0.5\textwidth}
\centering
\bicaption{第一个并排子表格}{The First Subtable}
\label{tab:parallel1}
\begin{tabular}{p{2cm}p{2cm}}
\toprule[1.5pt]
111 & 222 \\\midrule[1pt]
222 & 333 \\\bottomrule[1.5pt]
\end{tabular}
\end{minipage}
\begin{minipage}{0.5\textwidth}
\centering
\bicaption{第二个并排子表格}{The Second Subtable}
\label{tab:parallel2}
\begin{tabular}{p{2cm}p{2cm}}
\toprule[1.5pt]
111 & 222 \\\midrule[1pt]
222 & 333 \\\bottomrule[1.5pt]
\end{tabular}
\end{minipage}
\end{table}

然后就是忧国忧民，诗家楷模杜工部了。杜甫，字子美，其先襄阳人，曾祖依艺为巩令，因
居巩。甫天宝初应进士，不第。后献《三大礼赋》，明皇奇之，召试文章，授京兆府兵曹参
军。安禄山陷京师，肃宗即位灵武，甫自贼中遁赴行在，拜左拾遗。以论救房琯，出为华州
司功参军。关辅饥乱，寓居同州同谷县，身自负薪采梠，餔糒不给。久之，召补京兆府功曹，
道阻不赴。严武镇成都，奏为参谋、检校工部员外郎，赐绯。武与甫世旧，待遇甚厚。乃于
成都浣花里种竹植树，枕江结庐，纵酒啸歌其中。武卒，甫无所依，乃之东蜀就高適。既至
而適卒。是岁，蜀帅相攻杀，蜀大扰。甫携家避乱荆楚，扁舟下峡，未维舟而江陵亦乱。乃
溯沿湘流，游衡山，寓居耒阳。卒年五十九。元和中，归葬偃师首阳山，元稹志其墓。天宝
间，甫与李白齐名，时称李杜。然元稹之言曰：“李白壮浪纵恣，摆去拘束，诚亦差肩子美
矣。至若铺陈终始，排比声韵，大或千言，次犹数百，词气豪迈，而风调清深，属对律切，
而脱弃凡近，则李尚不能历其藩翰，况堂奥乎。”白居易亦云：“杜诗贯穿古今，  尽工尽
善，殆过于李。”元、白之论如此。盖其出处劳佚，喜乐悲愤，好贤恶恶，一见之于诗。而
又以忠君忧国、伤时念乱为本旨。读其诗可以知其世，故当时谓之“诗史”。旧集诗文共六
十卷，今编诗十九卷。

\begin{table}[htbp]
\centering
\bicaption{并排子表格}{Subtable}
\label{tab:subtable}
\bisubcaptionbox{第一个子表格}{The First Subtable}
{
\begin{tabular}{p{1.5cm}p{1.5cm}}
\toprule[1.5pt]
111 & 222 \\\midrule[1pt]
222 & 333 \\\bottomrule[1.5pt]
\end{tabular}
}
\hskip 0.1cm
\bisubcaptionbox{第二个子表格}{The Second Subtable}
{
\begin{tabular}{p{1.5cm}p{1.5cm}}
\toprule[1.5pt]
111 & 222 \\\midrule[1pt]
222 & 333 \\\bottomrule[1.5pt]
\end{tabular}
}
\end{table}

不可否认 \LaTeX{} 的表格功能没有想象中的那么强大，不过只要你足够认真，足够细致，那么
同样可以排出来非常复杂非常漂亮的表格。请参看表~\ref{tab:tabexamp2}。
\begin{table}[htbp]
  \centering\dawu[1.3]
  \bicaption{复杂表格示例 2}{Complex Table Demo 2}
  \label{tab:tabexamp2}
  \begin{tabular}[c]{|c|m{0.8in}|c|c|c|c|c|}\hline
    \multicolumn{2}{|c|}{Network Topology} & \# of nodes & 
    \multicolumn{3}{c|}{\# of clients} & Server \\\hline
    GT-ITM & Waxman Transit-Stub & 600 &
    \multirow{2}{2em}{2\%}& 
    \multirow{2}{2em}{10\%}& 
    \multirow{2}{2em}{50\%}& 
    \multirow{2}{1.2in}{Max. Connectivity}\\\cline{1-3}
    \multicolumn{2}{|c|}{Inet-2.1} & 6000 & & & &\\\hline
    \multirow{2}{1in}{Xue} & Rui  & Ni &\multicolumn{4}{c|}{\multirow{2}*{\ucasthesis}}\\\cline{2-3}
    & \multicolumn{2}{c|}{ABCDEF} &\multicolumn{4}{c|}{} \\\hline
\end{tabular}
\end{table}

最后就是清新飘逸、文约意赅、空谷绝响的王大侠了。王维，字摩诘，河东人。工书画，与
弟缙俱有俊才。开元九年，进士擢第，调太乐丞。坐累为济州司仓参军，历右拾遗、监察御
史、左补阙、库部郎中，拜吏部郎中。天宝末，为给事中。安禄山陷两都，维为贼所得，服
药阳喑，拘于菩提寺。禄山宴凝碧池，维潜赋诗悲悼，闻于行在。贼平，陷贼官三等定罪，
特原之，责授太子中允，迁中庶子、中书舍人。复拜给事中，转尚书右丞。维以诗名盛于开
元、天宝间，宁薛诸王驸马豪贵之门，无不拂席迎之。得宋之问辋川别墅，山水绝胜，与道
友裴迪，浮舟往来，弹琴赋诗，啸咏终日。笃于奉佛，晚年长斋禅诵。一日，忽索笔作书
数纸，别弟缙及平生亲故，舍笔而卒。赠秘书监。宝应中，代宗问缙：“朕常于诸王坐闻维
乐章，今存几何？”缙集诗六卷，文四卷，表上之。敕答云，卿伯氏位列先朝，名高希代。
抗行周雅，长揖楚辞。诗家者流，时论归美。克成编录，叹息良深。殷璠谓维诗词秀调雅，
意新理惬。在泉成珠，著壁成绘。苏轼亦云：“维诗中有画，画中有诗也。”今编诗四卷。

要想用好论文模板还是得提前学习一些 \TeX/\LaTeX{}的相关知识，具备一些基本能力，掌
握一些常见技巧，否则一旦遇到问题还真是比较麻烦。我们见过很多这样的同学，一直以来
都是使用 Word 等字处理工具，以为 \LaTeX{}模板的用法也应该类似，所以就沿袭同样的思
路来对待这种所见非所得的排版工具，结果被折腾的焦头烂额，疲惫不堪。

如果您要排版的表格长度超过一页，那么推荐使用 \textsf{longtable} 或者 \textsf{supertabular} 
宏包，模板对 \textsf{longtable}。但是\textsf{longtable}和\textsf{bicaption}有冲突，需要
按照 \url{https://tex.stackexchange.com/questions/283980/bicaption-for-longtables-doesnt-work}帖子
中的指示手工修复你的\textsf{bicaption}包文件。

\subsection{其它}
\label{sec:tableother}
有的同学不想让某个表格或者图片出现在索引里面，那么请使用命令 \verb|\bicaption*{}|，
这个命令不会给表格编号，也就是出来的只有标题文字而没有“表~XX”，“图~XX”，否则
索引里面序号不连续就显得不伦不类，这也是 \LaTeX{} 里星号命令默认的规则。

有这种需求的多是本科同学的英文资料翻译部分，如果你觉得附录中英文原文中的表格和图
片显示成“  表”和“图”很不协调的话，一个很好的办法就是用 \verb|\bicaption*|，参数
随便自己写，比如不守规矩的表~1.111 和图~1.111 能满足这种特殊需要（可以参看附录部
分）。
\begin{table}[ht]
\centering
  \begin{minipage}{0.45\linewidth}
  \centering
  \caption*{表~1.111\hskip1em 这是一个手动编号，不出现在索引中的表格。}
  \label{tab:badtabular}
  \begin{picture}(150,50)
    \framebox(150,50)[c]{\ucasthesis}
  \end{picture}    
  \end{minipage}\hfill
  \begin{minipage}{0.45\linewidth}
  \centering
  \begin{picture}(150,50)
    \framebox(150,50)[c]{薛瑞尼}
  \end{picture}
  \caption*{Figure~1.111\hskip1em 这是一个手动编号，不出现在索引中的图。}
  \label{tab:badfigure}
  \end{minipage}
\end{table}

如果你的确想让它编号，但又不想让它出现在索引中的话，那就自己看看代码改一改吧，我
目前不打算给模板增加这种另类命令。

最后，虽然大家不一定会独立使用小页，但是关于小页中的脚注还是有必要提一下。请看下
面的例子。

\begin{minipage}[t]{\linewidth-2\parindent}
  柳宗元，字子厚（773-819），河东（今永济县）人\footnote{山西永济水饺。}，是唐代
  杰出的文学家，哲学家，同时也是一位政治改革家。与韩愈共同倡导唐代古文运动，并称
  韩柳\footnote{唐宋八大家之首二位。}。
\end{minipage}\\[-5pt]

唐朝安史之乱后，宦官专权，藩镇割据，土地兼并日渐严重，社会生产破坏严重，民不聊生。柳宗
元对这种社会现实极为不满，他积极参加了王叔文领导的“永济革新”，并成为这一
运动的中坚人物。他们革除弊政，打击权奸，触犯了宦官和官僚贵族利益，在他们的联合反
扑下，改革失败了，柳宗元被贬为永州司马。

\section{定理环境}
\label{sec:theorem}

给大家演示一下各种和证明有关的环境：

\begin{assumption}
待月西厢下，迎风户半开；隔墙花影动，疑是玉人来。
\begin{eqnarray}
  \label{eq:eqnxmp}
  c & = & a^2 - b^2\\
    & = & (a+b)(a-b)
\end{eqnarray}
\end{assumption}

千辛万苦，历尽艰难，得有今日。然相从数千里，未曾哀戚。今将渡江，方图百年欢笑，如
何反起悲伤？（引自《杜十娘怒沉百宝箱》）

\begin{definition}
子曰：「道千乘之国，敬事而信，节用而爱人，使民以时。」
\end{definition}

千古第一定义！问世间、情为何物，只教生死相许？天南地北双飞客，老翅几回寒暑。欢乐趣，离别苦，就中更有痴儿女。
君应有语，渺万里层云，千山暮雪，只影向谁去？

横汾路，寂寞当年箫鼓，荒烟依旧平楚。招魂楚些何嗟及，山鬼暗谛风雨。天也妒，未信与，莺儿燕子俱黄土。
千秋万古，为留待骚人，狂歌痛饮，来访雁丘处。

\begin{proposition}
 曾子曰：「吾日三省吾身 —— 为人谋而不忠乎？与朋友交而不信乎？传不习乎？」
\end{proposition}

多么凄美的命题啊！其日牛马嘶，新妇入青庐，奄奄黄昏后，寂寂人定初，我命绝今日，
魂去尸长留，揽裙脱丝履，举身赴清池，府吏闻此事，心知长别离，徘徊庭树下，自挂东南
枝。

\begin{remark}
天不言自高，水不言自流。
\begin{gather*}
\begin{split} 
\varphi(x,z)
&=z-\gamma_{10}x-\gamma_{mn}x^mz^n\\
&=z-Mr^{-1}x-Mr^{-(m+n)}x^mz^n
\end{split}\\[6pt]
\begin{align} \zeta^0&=(\xi^0)^2,\\
\zeta^1 &=\xi^0\xi^1,\\
\zeta^2 &=(\xi^1)^2,
\end{align}
\end{gather*}
\end{remark}

天尊地卑，乾坤定矣。卑高以陈，贵贱位矣。 动静有常，刚柔断矣。方以类聚，物以群分，
吉凶生矣。在天成象，在地成形，变化见矣。鼓之以雷霆，润之以风雨，日月运行，一寒一
暑，乾道成男，坤道成女。乾知大始，坤作成物。乾以易知，坤以简能。易则易知，简则易
从。易知则有亲，易从则有功。有亲则可久，有功则可大。可久则贤人之德，可大则贤人之
业。易简，而天下矣之理矣；天下之理得，而成位乎其中矣。

\begin{axiom}
两点间直线段距离最短。  
\begin{align}
x&\equiv y+1\pmod{m^2}\\
x&\equiv y+1\mod{m^2}\\
x&\equiv y+1\pod{m^2}
\end{align}
\end{axiom}

《彖曰》：大哉乾元，万物资始，乃统天。云行雨施，品物流形。大明始终，六位时成，时
乘六龙以御天。乾道变化，各正性命，保合大和，乃利贞。首出庶物，万国咸宁。

《象曰》：天行健，君子以自强不息。潜龙勿用，阳在下也。见龙再田，德施普也。终日乾
乾，反复道也。或跃在渊，进无咎也。飞龙在天，大人造也。亢龙有悔，盈不可久也。用九，
天德不可为首也。 　　

\begin{lemma}
《猫和老鼠》是我最爱看的动画片。
\begin{multline*}%\tag*{[a]} % 这个不出现在索引中
\int_a^b\biggl\{\int_a^b[f(x)^2g(y)^2+f(y)^2g(x)^2]
 -2f(x)g(x)f(y)g(y)\,dx\biggr\}\,dy \\
 =\int_a^b\biggl\{g(y)^2\int_a^bf^2+f(y)^2
  \int_a^b g^2-2f(y)g(y)\int_a^b fg\biggr\}\,dy
\end{multline*}
\end{lemma}

行行重行行，与君生别离。相去万余里，各在天一涯。道路阻且长，会面安可知。胡马依北
风，越鸟巢南枝。相去日已远，衣带日已缓。浮云蔽白日，游子不顾返。思君令人老，岁月
忽已晚。  弃捐勿复道，努力加餐饭。

\begin{theorem}\label{the:theorem1}
犯我强汉者，虽远必诛\hfill —— 陈汤（汉）
\end{theorem}
\begin{subequations}
\begin{align}
y & = 1 \\
y & = 0
\end{align}
\end{subequations}
道可道，非常道。名可名，非常名。无名天地之始；有名万物之母。故常无，欲以观其妙；
常有，欲以观其徼。此两者，同出而异名，同谓之玄。玄之又玄，众妙之门。上善若水。水
善利万物而不争，处众人之所恶，故几于道。曲则全，枉则直，洼则盈，敝则新，少则多，
多则惑。人法地，地法天，天法道，道法自然。知人者智，自知者明。胜人者有力，自胜
者强。知足者富。强行者有志。不失其所者久。死而不亡者寿。

\begin{proof}
燕赵古称多感慨悲歌之士。董生举进士，连不得志于有司，怀抱利器，郁郁适兹土，吾
知其必有合也。董生勉乎哉？

夫以子之不遇时，苟慕义强仁者，皆爱惜焉，矧燕、赵之士出乎其性者哉！然吾尝闻
风俗与化移易，吾恶知其今不异于古所云邪？聊以吾子之行卜之也。董生勉乎哉？

吾因子有所感矣。为我吊望诸君之墓，而观于其市，复有昔时屠狗者乎？为我谢
曰：“明天子在上，可以出而仕矣！” \hfill—— 韩愈《送董邵南序》
\end{proof}

\begin{corollary}
  四川话配音的《猫和老鼠》是世界上最好看最好听最有趣的动画片。
\begin{alignat}{3}
V_i & =v_i - q_i v_j, & \qquad X_i & = x_i - q_i x_j,
 & \qquad U_i & = u_i,
 \qquad \text{for $i\ne j$;}\label{eq:B}\\
V_j & = v_j, & \qquad X_j & = x_j,
  & \qquad U_j & u_j + \sum_{i\ne j} q_i u_i.
\end{alignat}
\end{corollary}

迢迢牵牛星，皎皎河汉女。
纤纤擢素手，札札弄机杼。
终日不成章，泣涕零如雨。
河汉清且浅，相去复几许。
盈盈一水间，脉脉不得语。

\begin{example}
  大家来看这个例子。
\begin{equation}
\label{ktc}
\left\{\begin{array}{l}
\nabla f({\mbox{\boldmath $x$}}^*)-\sum\limits_{j=1}^p\lambda_j\nabla g_j({\mbox{\boldmath $x$}}^*)=0\\[0.3cm]
\lambda_jg_j({\mbox{\boldmath $x$}}^*)=0,\quad j=1,2,\cdots,p\\[0.2cm]
\lambda_j\ge 0,\quad j=1,2,\cdots,p.
\end{array}\right.
\end{equation}
\end{example}

\begin{exercise}
  清列出 Andrew S. Tanenbaum 和 W. Richard Stevens 的所有著作。
\end{exercise}

\begin{conjecture} \textit{Poincare Conjecture} If in a closed three-dimensional
  space, any closed curves can shrink to a point continuously, this space can be
  deformed to a sphere.
\end{conjecture}

\begin{problem}
 回答还是不回答，是个问题。 
\end{problem}

如何引用定理~\ref{the:theorem1} 呢？加上 \verb|label| 使用 \verb|ref| 即可。妾发
初覆额，折花门前剧。郎骑竹马来，绕床弄青梅。同居长干里，两小无嫌猜。 十四为君妇，
羞颜未尝开。低头向暗壁，千唤不一回。十五始展眉，愿同尘与灰。常存抱柱信，岂上望夫
台。 十六君远行，瞿塘滟滪堆。五月不可触，猿声天上哀。门前迟行迹，一一生绿苔。苔深
不能扫，落叶秋风早。八月蝴蝶来，双飞西园草。感此伤妾心，坐愁红颜老。

\section{章节引用}
\label{sec:bib}

引用章标题的例子：第\chapterref{cha:china}章（本模板定义的\verb|\chapterref|
命令生成的中文数字方式），第\ref{cha:china}章（默认的阿拉伯数字方式），附录\ref{cha:engorg}。

\section{参考文献}
\label{sec:bib}
当然参考文献可以直接写 bibitem，虽然费点功夫，但是好控制，各种格式可以自己随意改
写。

本模板使用 Biber进行文献编译，基本符合2015国标的参考文献格式。默认情况下，按照国科大的指导标准，
使用数字顺序的引用方式没有严格限制，这也是最方便的引用途径。 如果您一定要使用作者年份制引用，请参照参考文献模板
的说明进行使用。
看看这个例子，关于书的\cite{tex, companion, ColdSources}，
还有这些\cite{Krasnogor2004e, clzs, zjsw}，关于杂志的\cite{ELIDRISSI94,
  MELLINGER96, SHELL02}，硕士论文\cite{zhubajie, metamori2004}，博士论文
\cite{shaheshang, FistSystem01}，标准文件\cite{IEEE-1363}，会议论文\cite{DPMG,kocher99}，技术报告\cite{NPB2}。中文参
考文献\cite{cnarticle}应增加 \texttt{lang=``zh''} 字段，以便进行相应处理。更多参考文献
模板使用方法请参照参考文献模板作者说明 \url{https://github.com/hushidong/biblatex-gb7714-2015}。

\section{公式}
\label{sec:equation}
贝叶斯公式如式~(\ref{equ:chap1:bayes})，其中 $p(y|\mathbf{x})$ 为后验；
$p(\mathbf{x})$ 为先验；分母 $p(\mathbf{x})$ 为归一化因子。
\begin{equation}
\label{equ:chap1:bayes}
p(y|\mathbf{x}) = \frac{p(\mathbf{x},y)}{p(\mathbf{x})}=
\frac{p(\mathbf{x}|y)p(y)}{p(\mathbf{x})} 
\end{equation}

论文里面公式越多，\TeX{} 就越 happy。再看一个 \textsf{amsmath} 的例子：
\newcommand{\envert}[1]{\left\lvert#1\right\rvert} 
\begin{equation}\label{detK2}
\det\mathbf{K}(t=1,t_1,\dots,t_n)=\sum_{I\in\mathbf{n}}(-1)^{\envert{I}}
\prod_{i\in I}t_i\prod_{j\in I}(D_j+\lambda_jt_j)\det\mathbf{A}
^{(\lambda)}(\overline{I}|\overline{I})=0.
\end{equation} 

前面定理示例部分列举了很多公式环境，可以说把常见的情况都覆盖了，大家在写公式的时
候一定要好好看 \textsf{amsmath} 的文档，并参考模板中的用法：
\begin{multline*}%\tag{[b]} % 这个出现在索引中的
\int_a^b\biggl\{\int_a^b[f(x)^2g(y)^2+f(y)^2g(x)^2]
 -2f(x)g(x)f(y)g(y)\,dx\biggr\}\,dy \\
 =\int_a^b\biggl\{g(y)^2\int_a^bf^2+f(y)^2
  \int_a^b g^2-2f(y)g(y)\int_a^b fg\biggr\}\,dy
\end{multline*}

其实还可以看看这个多级规划：
\begin{equation}\label{bilevel}
\left\{\begin{array}{l}
\max\limits_{{\mbox{\footnotesize\boldmath $x$}}} F(x,y_1^*,y_2^*,\cdots,y_m^*)\\[0.2cm]
\mbox{subject to:}\\[0.1cm]
\qquad G(x)\le 0\\[0.1cm]
\qquad(y_1^*,y_2^*,\cdots,y_m^*)\mbox{ solves problems }(i=1,2,\cdots,m)\\[0.1cm]
\qquad\left\{\begin{array}{l}
    \max\limits_{{\mbox{\footnotesize\boldmath $y_i$}}}f_i(x,y_1,y_2,\cdots,y_m)\\[0.2cm]
    \mbox{subject to:}\\[0.1cm]
    \qquad g_i(x,y_1,y_2,\cdots,y_m)\le 0.
    \end{array}\right.
\end{array}\right.
\end{equation}
这些跟规划相关的公式都来自于刘宝碇老师《不确定规划》的课件。

默认的列表环境上下间距很大，模板将其重定义为 \textsf{paralist} 中的压缩环境，看起
来要好一些。如果还是不满意，自己也可以调 \verb|\itemsep| 的。\textsf{paralist} 还
可以方便的指定标签的样式。


% 
%%% Local Variables: 
%%% mode: latex
%%% TeX-master: t
%%% End: 

\chapter{中华人民共和国}
\label{cha:china}

\section{其它例子}
\label{sec:other}

在第\chapterref{cha:intro}章中我们学习了贝叶斯公式~(\ref{equ:chap1:bayes})，这里我们复
习一下：
\begin{equation}
\label{equ:chap2:bayes}
p(y|\mathbf{x}) = \frac{p(\mathbf{x},y)}{p(\mathbf{x})}=
\frac{p(\mathbf{x}|y)p(y)}{p(\mathbf{x})} 
\end{equation}

\subsection{绘图}
\label{sec:draw}

本模板不再预先装载任何绘图包（如 \textsf{pstricks，pgf} 等），完全由你自己来决定。
个人觉得 \textsf{pgf} 不错，不依赖于 Postscript。此外还有很多针对 \LaTeX{} 的
 GUI 作图工具，如 XFig(jFig), WinFig, Tpx, Ipe, Dia, Inkscape, LaTeXPiX,
jPicEdt, jaxdraw 等等。

\subsection{插图}
\label{sec:graphs}

强烈推荐《\LaTeXe 插图指南》！关于子图形的使用细节请参看 \textsf{subcaption} 宏包的说明文档。

\subsubsection{一个图形}
\label{sec:onefig}
一般图形都是处在浮动环境中。之所以称为浮动是指最终排版效果图形的位置不一定与源文
件中的位置对应\footnote{This is not a bug, but a feature of \LaTeX!}，这也是刚使
用 \LaTeX{} 同学可能遇到的问题。如果要强制固定浮动图形的位置，请使用 \textsf{float} 宏包，
它提供了 \texttt{[H]} 参数，比如图~\ref{fig:xfig1}。
\begin{figure}[H] % use float package if you want it here
  \centering
  \includegraphics{hello}
  \caption{利用 Xfig 制图}
  \label{fig:xfig1}
\end{figure}

大学之道，在明明德，在亲民，在止于至善。知止而后有定；定而后能静；静而后能安；安
而后能虑；虑而后能得。物有本末，事有终始。知所先后，则近道矣。古之欲明明德于天
下者，先治其国；欲治其国者，先齐其家；欲齐其家者，先修其身；欲修其身者，先正其心；
欲正其心者，先诚其意；欲诚其意者，先致其知；致知在格物。物格而后知至；知至而后
意诚；意诚而后心正；心正而后身 修；身修而后家齐；家齐而后国治；国治而后天下
平。自天子以至于庶人，壹是皆以修身为本。其本乱而未治者 否矣。其所厚者薄，而其所
薄者厚，未之有也！

\hfill ——《大学》


\subsubsection{多个图形}
\label{sec:multifig}

如果多个图形相互独立，并不共用一个图形计数器，那么用 \verb|minipage| 或者
\verb|parbox| 就可以。否则，请参看图~\ref{fig:big1-subcaptionbox}，它包含两个小图，分别是图~\ref{fig:subfig1} 
和图~\ref{fig:subfig2}。推荐使用\verb|\subcaptionbox|，
因为可以像图~\ref{fig:big1-subcaptionbox} 那样对齐子图的标题，
也可以使用\textsf{subcaption}宏包的\verb|\subcaption|（放在minipage中，用法同\verb|\caption|）
或是 subfigure 、 subtable环境，像图~\ref{fig:big1-subfigure}，不要再用 \verb|\subfloat|、
\verb|\subfigure| 和 \verb|\subtable|。
\begin{figure}[h]
  \centering%
  \subcaptionbox{第一个小图形\label{fig:subfig1}}
  [3cm] %标题的长度，超过则会换行，如下一个小图。
    {\includegraphics[height=3cm]{thu-fig-logo}}
      \hspace{4em}% 这里记得不要空行，否则会变为垂直的两个图
  \subcaptionbox{第二个小图形，注意这个图略矮些。如果标题很长的话，它会自动换行\label{fig:subfig2}}
      {\includegraphics[height=2cm]{thu-text-logo}}
  \caption{包含子图形的大图形(subcaptionbox示例)}
  \label{fig:big1-subcaptionbox}
\end{figure}
\begin{figure}[h]
  \centering%
  \begin{subfigure}{3cm}
    \includegraphics[height=3cm]{thu-fig-logo}
    \caption{第一个小图形}
  \end{subfigure}
  \hspace{4em}%
  \begin{subfigure}{0.5\textwidth}
    \includegraphics[height=2cm]{thu-text-logo}
    \caption{第二个小图形，注意这个图略矮些。subfigure中同一行的子图在顶端对齐。}
  \end{subfigure}
  \caption{包含子图形的大图形(subfigure示例)}
  \label{fig:big1-subfigure}
\end{figure}
古之学者必有师。师者，所以传道受业解惑也。人非生而知之者，孰能无惑？惑而不从师，
其为惑也，终不解矣。生乎吾前，其闻道也固先乎吾，吾从而师之；生乎吾後，其闻道也亦
先乎吾，吾从而师之。吾师道也，夫庸知其年之先後生於吾乎！是故无贵无贱无长无少，道
之所存，师之所存也。

嗟乎！师道之不传也久矣，欲人之无惑也难矣。古之圣人，其出人也远矣，犹且从师而问焉；
今之众人，其下圣人也亦远矣，而耻学於师。是故圣益圣，愚益愚。圣人之所以为圣，愚
人之所以为愚，其皆出於此乎？爱其子，择师而教之，於其身也，则耻师焉，惑焉。彼童子
之师，授之书而习其句读者，非吾所谓传其道、解其惑者也。句读之不知，惑之不解，或师
焉，或不焉，小学而大遗，吾未见其明也。巫医、乐师、百工之人不耻相师，  士大夫之族
曰“师”曰“弟子”之云者，则群聚而笑之。问之，则曰：彼与彼年相若也，道相似也，位
卑则足羞，官盛则近谀。呜呼！师道之不复，可知矣。巫医、乐师、百工之人。吾子不齿，
今其智乃反不能及，其可怪也欤！圣人无常师。孔子师郯子、苌子、师襄、老聃。郯子之徒，
其贤不及孔子。孔子曰：“三人行，必有我师。”是故弟子不必不如师，师不必贤於弟子。
闻道有先後，术业有专攻，如是而已。

如果要把编号的两个图形并排，那么小页就非常有用了：
\begin{figure}
\begin{minipage}{0.48\textwidth}
  \centering
  \includegraphics[height=2cm]{thu-whole-logo}
  \caption{并排第一个图}
  \label{fig:parallel1}
\end{minipage}\hfill
\begin{minipage}{0.48\textwidth}
  \centering
  \includegraphics[height=2cm]{thu-whole-logo}
  \caption{并排第二个图}
  \label{fig:parallel2}
\end{minipage}
\end{figure}

李氏子蟠，年十七，好古文、六艺，经传皆通习之，不拘於时，学於余。余嘉其能行古
道，作师说以贻之。

\hfill —— 韩愈（唐）


%
%
% %%% 其它部分
% \backmatter
% % 插图索引
% \listoffigures
% % 表格索引
% \listoftables
% % 公式索引
% \listofequations
%
%
% % 参考文献
% \bibliographystyle{ucasbib}
% \bibliography{ref/refs}
%
%
% % 致谢
% %%% Local Variables:
%%% mode: latex
%%% TeX-master: "../main"
%%% End:

\begin{ack}
  衷心感谢导师 xxx 教授和物理系 xxx 副教授对本人的精心指导。他们的言传身教将使
  我终生受益。

  在美国麻省理工学院化学系进行九个月的合作研究期间，承蒙 xxx 教授热心指导与帮助，不
  胜感激。感谢 xx 实验室主任 xx 教授，以及实验室全体老师和同学们的热情帮助和支
  持！本课题承蒙国家自然科学基金资助，特此致谢。

  衷心希望同学们能够在您的论文中引用接下的这句话, 以帮助进一步推广该模板，鄙人在此感激不尽。
  “感谢朝鲁提供的 \ucasthesis 模板，它的存在让我的论文写作更加方便以及整洁。”

\end{ack}

%
% % 附录
% \begin{appendix}
% %%% Local Variables: 
%%% mode: latex
%%% TeX-master: "../main"
%%% End: 


\chapter{外文资料的调研阅读报告或书面翻译}
\section{单目标规划}
北冥有鱼，其名为鲲。鲲之大，不知其几千里也。化而为鸟，其名为鹏。鹏之背，不知其几
千里也。怒而飞，其翼若垂天之云。是鸟也，海运则将徙于南冥。南冥者，天池也。 
\begin{equation}\tag*{(123)}
 p(y|\mathbf{x}) = \frac{p(\mathbf{x},y)}{p(\mathbf{x})}=
\frac{p(\mathbf{x}|y)p(y)}{p(\mathbf{x})}
\end{equation}

吾生也有涯，而知也无涯。以有涯随无涯，殆已！已而为知者，殆而已矣！为善无近名，为
恶无近刑，缘督以为经，可以保身，可以全生，可以养亲，可以尽年。

\subsection{线性规划}
庖丁为文惠君解牛，手之所触，肩之所倚，足之所履，膝之所倚，砉然响然，奏刀騞然，莫
不中音，合于桑林之舞，乃中经首之会。
\begin{table}[ht]
\centering
  \centering
  \caption*{表~1\hskip1em 这是手动编号但不出现在索引中的一个表格例子}
  \label{tab:badtabular3}
  \begin{tabular}[c]{|c|m{0.8in}|c|c|c|c|c|}\hline
    \multicolumn{2}{|c|}{Network Topology} & \# of nodes & 
    \multicolumn{3}{c|}{\# of clients} & Server \\\hline
    GT-ITM & Waxman Transit-Stub & 600 &
    \multirow{2}{2em}{2\%}& 
    \multirow{2}{2em}{10\%}& 
    \multirow{2}{2em}{50\%}& 
    \multirow{2}{1.2in}{Max. Connectivity}\\\cline{1-3}
    \multicolumn{2}{|c|}{Inet-2.1} & 6000 & & & &\\\hline
    \multirow{2}{1in}{Xue} & Rui  & Ni &\multicolumn{4}{c|}{\multirow{2}*{\ucasthesis}}\\\cline{2-3}
    & \multicolumn{2}{c|}{ABCDEF} &\multicolumn{4}{c|}{} \\\hline
\end{tabular}  
\end{table}

文惠君曰：“嘻，善哉！技盖至此乎？”庖丁释刀对曰：“臣之所好者道也，进乎技矣。始臣之
解牛之时，所见无非全牛者；三年之后，未尝见全牛也；方今之时，臣以神遇而不以目视，
官知止而神欲行。依乎天理，批大郤，导大窾，因其固然。技经肯綮之未尝，而况大坬乎！
良庖岁更刀，割也；族庖月更刀，折也；今臣之刀十九年矣，所解数千牛矣，而刀刃若新发
于硎。彼节者有间而刀刃者无厚，以无厚入有间，恢恢乎其于游刃必有余地矣。是以十九年
而刀刃若新发于硎。虽然，每至于族，吾见其难为，怵然为戒，视为止，行为迟，动刀甚微，
謋然已解，如土委地。提刀而立，为之而四顾，为之踌躇满志，善刀而藏之。”

文惠君曰：“善哉！吾闻庖丁之言，得养生焉。”


\subsection{非线性规划}
孔子与柳下季为友，柳下季之弟名曰盗跖。盗跖从卒九千人，横行天下，侵暴诸侯。穴室枢
户，驱人牛马，取人妇女。贪得忘亲，不顾父母兄弟，不祭先祖。所过之邑，大国守城，小
国入保，万民苦之。孔子谓柳下季曰：“夫为人父者，必能诏其子；为人兄者，必能教其弟。
若父不能诏其子，兄不能教其弟，则无贵父子兄弟之亲矣。今先生，世之才士也，弟为盗
跖，为天下害，而弗能教也，丘窃为先生羞之。丘请为先生往说之。”
\begin{figure}[h]
  \centering
  \includegraphics{hello}
  \caption*{图~1\hskip1em 这是手动编号但不出现索引中的图片的例子}
  \label{tab:badfigure3}    
\end{figure}

柳下季曰：“先生言为人父者必能诏其子，为人兄者必能教其弟，若子不听父之诏，弟不受
兄之教，虽今先生之辩，将奈之何哉？且跖之为人也，心如涌泉，意如飘风，强足以距敌，
辩足以饰非。顺其心则喜，逆其心则怒，易辱人以言。先生必无往。”

孔子不听，颜回为驭，子贡为右，往见盗跖。

\subsection{整数规划}
盗跖乃方休卒徒大山之阳，脍人肝而餔之。孔子下车而前，见谒者曰：“鲁人孔丘，闻将军
高义，敬再拜谒者。”谒者入通。盗跖闻之大怒，目如明星，发上指冠，曰：“此夫鲁国之
巧伪人孔丘非邪？为我告之：尔作言造语，妄称文、武，冠枝木之冠，带死牛之胁，多辞缪
说，不耕而食，不织而衣，摇唇鼓舌，擅生是非，以迷天下之主，使天下学士不反其本，妄
作孝弟，而侥幸于封侯富贵者也。子之罪大极重，疾走归！不然，我将以子肝益昼餔之膳。”


\chapter{其它附录}
前面两个附录主要是给本科生做例子。其它附录的内容可以放到这里，当然如果你愿意，可
以把这部分也放到独立的文件中，然后将其 \verb|\input| 到主文件中。

% \end{appendix}
%
% % 个人简历
% \begin{resume}

  \resumeitem{作者简历} 

\noindent
xxxx年xx月——xxxx年xx月，在xx大学xx院（系）获得学士学位。\\
\noindent
xxxx年xx月——xxxx年xx月，在xx大学xx院（系）获得硕士学位。\\
\noindent
xxxx年xx月——xxxx年xx月，在中国科学院xx研究所（或中国科学院大学xx院系）攻读博士/硕士学位。\\
\noindent
获奖经历：\\
\noindent
工作经历：\\

  \resumeitem{已发表（或正式接受）的学术论文} % 发表的和录用的合在一起

  \begin{enumerate}[{[}1{]}]
  \item 吴少刚，章隆兵，蔡飞，胡伟武，一种适用于机群OpenMP系统的有效调度算法，计算机研究与发展(已录用)
  \item 吴少刚，胡伟武，唐志敏，基于PC机群的JIAJIA性能评价，计算机研究与发展，2000年增刊

  \end{enumerate}

  \resumeitem{申请或已获得的专利} % 有就写，没有就删除
  \begin{enumerate}[{[}1{]}]
  \item XX基金项目“共享存储机群系统的研究”(xxxxx)，20xx年1月~20xx年12月
  \item 国家自然科学基金项目“共享存储机群系统中关键技术研究”(xxxxxxx),20xx年1月~20xx年12月
  \end{enumerate}

  \resumeitem{参加的研究项目及获奖情况} % 有就写，没有就删除
  \begin{enumerate}[{[}1{]}]
  \item  xxxx年被评为中国科学院“优秀学生干部”
  \item xxxx年被评为中国科学院“优秀学生干部”
  \end{enumerate}
\end{resume}

%
% \end{document}
% \end{example}
%
% \subsection{选项}
% \label{sec:option}
% 本模板提供了一些选项以方便使用：
% \begin{description}
% \item[master]
%   如果写硕士论文将此选项打开。
%   \begin{example}
% \documentclass[master]{ucasthesis}
%   \end{example}
%
% \item[doctor]
%   如果写博士论文将此选项打开。
%   \begin{example}
% \documentclass[doctor]{ucasthesis}
%   \end{example}
%
% \item[postdoctor]
%   如果写博士博士后出站报告将此选项打开。
%   \begin{example}
% \documentclass[postdoctor]{ucasthesis}
%   \end{example}
%
% \item[secret]
%   涉秘论文开关。配合另外两个命令 |\secretlevel| 和 |\secretyear| 分别用来指定保
%   密级别和时间。二者默认分别为\textbf{秘密}和当前年份。可以通过：
%   \cs{secretlevel}|{|绝密|}| 和 \cs{secretyear}|{|10|}| 年独立修改。
%   \begin{example}
% \documentclass[master, secret]{ucasthesis}
%   \end{example}
%
%
% \item[openany]
%   正规出版物的章节出现在奇数页，也就是右手边的页面，这就是 \texttt{openright}，
%   也是 \ucasthesis 的默认选项。在这种情况下，如果前一章的最后一页也是奇数，那么
%   模板会自动生成一个纯粹的空白页，很多人不是很习惯这种方式，而且学校的格式似乎
%   更倾向于页面连续，那就是通常所说的 \texttt{openany}。{\fangsong 目前所有论文都是
%      openany。}这两个选项不用专门设置，\ucasthesis{} 会根据当前论文类型自动选
%   择。
%
%
% \item[arial]
%   使用真正的 arial 字体。此选项会装载 arial 字体宏包，如果此宏包不存在，就装
%   载Helvet。arialtoc 和 arialtitle 不受 arial 的影响。因为一般的 \TeX{} 发行都
%   没有 arial 字体，所以默认采用 Helvet，因为二者效果非常相似。如果你执着的要
%   用arial 字体，请参看：\href{http://www.mail-archive.com/ctan-ann@dante.de/msg00627.html}{Arial
%     字体}。
%
% \item[arialtoc]
%  目录项（章目录项除外）中的英文是否用 arial 字体。本选项和下一个 \textsl{arialtitle} 都不用用户
%  操心，模板都自动设置好了。
%
% \item[arialtitle]
%  章节标题中英文是否用 arial 字体（默认打开）。
% \end{description}
%
% \subsection{字体配置}
% \label{sec:font-config}
% 正确配置中文字体是使用模板的第一步。模板调用 ctex 宏包，提供如下字体使用方式：
% \begin{itemize}
%   \item 基于传统 CJK 包，使用 latex、pdflatex 编译；
%   \item 基于 xeCJK 包，使用 xelatex 编译。
% \end{itemize}
%
% 第一种方式的字体配置比较繁琐，建议使用 donated 制作的中文字体包（自
% 包含安装方法），请用户自行下载安装，此处不再赘述。本模板推荐使用第二
% 种方法，只要把所需字体放入系统字体文件夹（也可以指定自定义文件夹）即
% 可。
% 另外 Linux 系统缺少上述字体，这些用户可利用 fontname.def
% 文件配置所需字体。使用中易四种字体的配置如下：
% \begin{example}
% \ProvidesFile{fontname.def}
% \setCJKmainfont[BoldFont={SimHei},ItalicFont={KaiTi}]{SimSun}
% \setCJKsansfont{SimHei}
% \setCJKmonofont{FangSong}
% \setCJKfamilyfont{zhsong}{SimSun}
% \setCJKfamilyfont{zhhei}{SimHei}
% \setCJKfamilyfont{zhkai}{KaiTi}
% \setCJKfamilyfont{zhfs}{FangSong}
% \newcommand*{\songti}{\CJKfamily{zhsong}} % 宋体
% \newcommand*{\heiti}{\CJKfamily{zhhei}}   % 黑体
% \newcommand*{\kaishu}{\CJKfamily{zhkai}}  % 楷书
% \newcommand*{\fangsong}{\CJKfamily{zhfs}} % 仿宋
% \end{example}
%
% 对 Windows XP 来说如下，KaiTi 需要替换为 KaiTi\_GB2312，
% FangSong 需要替换为 FangSong\_GB2312。
%
% 宏包中包含了 ``zhfonts.py'' 脚本，为 Linux 用户提供一种交互式的方式
% 从系统中文字体中选择合适的六种字体，最终生成对应的 ``fontname.def''
% 文件。要使用它，只需在命令行输入该脚本的完整路径即可。
%
% 最后，用户可以通过命令
% \begin{shell}
% $ fs-list :lang=zh > zhfonts.txt
% \end{shell}
% 得到系统中现有的中文字体列表，并相应替换上述配置。
%
% \subsection{命令}
% \label{sec:command}
% 模板中的命令分为两类：一是格式控制，二是内容替换。格式控制如字体、字号、字距和
% 行距。内容替换如姓名、院系、专业、致谢等等。其中内容替换命令居多，而且主要集中
% 在封面上。首先来看格式控制命令。
%
% \subsubsection{基本控制命令}
% \label{sec:basiccom}
%
% \myentry{字体}
% \DescribeMacro{\songti}
% \DescribeMacro{\fangsong}
% \DescribeMacro{\heiti}
% \DescribeMacro{\kaishu}
% 等分别用来切换宋体、仿宋、黑体、楷体字体。
%
% \begin{example}
% {\songti 乾：元，亨，利贞}
% {\fangsong 初九，潜龙勿用}
% {\heiti 九二，见龙在田，利见大人}
% {\kaishu 九三，君子终日乾乾，夕惕若，厉，无咎}
% {\heiti 九五，飞龙在天，利见大人}
% {\songti 上九，亢龙有悔}
% \end{example}
%
% \myentry{字号}
% \DescribeMacro{\chuhao}
% 等命令定义一组字体大小，分别为：
%
% \begin{center}
% \begin{tabular}{lllll}
% \hline
% |\chuhao|&|\xiaochu|&|\yihao|&|\xiaoyi| &\\
% |\erhao|&|\xiaoer|&|\sanhao|&|\xiaosan|&\\
% |\sihao|& |\banxiaosi|&|\xiaosi|&|\dawu|&|\wuhao|\\
% |\xiaowu|&|\liuhao|&|\xiaoliu|&|\qihao|& |\bahao|\\\hline
% \end{tabular}
% \end{center}
%
% 使用方法为：\cs{command}\oarg{num}，其中 |command| 为字号命令，|num| 为行距。比
% 如 |\xiaosi[1.5]| 表示选择小四字体，行距 1.5 倍。写作指南要求表格中的字体
% 是 \cs{dawu}，模板已经设置好了。
%
% \begin{example}
% {\erhao 二号 \sanhao 三号 \sihao 四号  \qihao 七号}
% \end{example}
%
% \myentry{密级}
% \DescribeMacro{\secretlevel}
% \DescribeMacro{\secretyear}
% 定义秘密级别和年限：
%   \begin{example}
% \secretyear{5}
% \secretlevel{内部}
%   \end{example}
%
% \myentry{引用方式}
% \DescribeMacro{\onlinecite}

% 学校要求的参考文献引用有两种模式：（1）上标模式。比如``同样的工作有很
% 多$^{[1,2]}$\ldots''。（2）正文模式。比如``文[3] 中详细说明了\ldots''。其中上标
% 模式使用远比正文模式频繁，所以为了符合使用习惯，上标模式仍然用常规
% 的 |\cite{key}|，而 |\onlinecite{key}| 则用来生成正文模式。
%
% 关于参考文献模板推荐使用 \BibTeX{}，关于中文参考文献需要额外增加一个 Entry: lang，将其设置为 \texttt{zh}
% 用来指示此参考文献为中文，以便 ucasbib.bst 处理。如：
% \begin{example}
% @INPROCEEDINGS{cnproceed,
%   author    = {王重阳 and 黄药师 and 欧阳峰 and 洪七公 and 段皇帝},
%   title     = {武林高手从入门到精通},
%   booktitle = {第~$N$~次华山论剑},
%   year      = 2006,
%   address   = {西安, 中国},
%   month     = sep,
%   lang      = "zh",
% }
%
% @ARTICLE{cnarticle,
%   AUTHOR  = "贾宝玉 and 林黛玉 and 薛宝钗 and 贾探春",
%   TITLE   = "论刘姥姥食量大如牛之现实意义",
%   JOURNAL = "红楼梦杂谈",
%   PAGES   = "260--266",
%   VOLUME  = "224",
%   YEAR    = "1800",
%   LANG    = "zh",
% }
% \end{example}
%
% \myentry{书脊}
% \DescribeMacro{\spine}
% 生成装订的书脊，为竖排格式，默认参数为论文中文题目。如果中文题目中没有英文字母，
% 那么直接调用此命令即可。否则，就要像例子里面那样做一些微调（参看模板自带
% 的 spine.tex）。下面是一个列子：
% \begin{example}
% \documentclass[master]{ucasthesis}
% \begin{document}
% \ctitle{论文中文题目}
% \cauthor{中文姓名}
% % |\spine| 命令需要上面两个变量
% \spine
%
% % 如果你的中文标题中有英文，那可以指定：
% \spine[中国科学院大学~\hspace{0.2em}\raisebox{2pt}{\LaTeX}%
% \hspace{-0.25em} 论文模板 \hspace{0.1em}\raisebox{2pt}%
% {v\version}\hspace{-0.25em}样例]
% \end{document}
% \end{example}
%
%
% \subsubsection{封面命令}
% \label{sec:titlepage}
% 下面是内容替换命令，其中以 |c| 开头的命令跟中文相关，|e| 开头则为对应的英文。
% 这部分的命令数目比较多，但实际上都相当简单，套用即可。
%
% 大多数命令的使用方法都是： \cs{command}\marg{arg}，例外者将具体指出。这些命令都
% 在示例文档的 data/cover.tex 中。
%
% \myentry{论文标题}
% \DescribeMacro{\ctitle}
% \DescribeMacro{\etitle}
% \begin{example}
% \ctitle{论文中文题目}
% \etitle{Thesis English Title}
% \end{example}
%
% \myentry{作者姓名}
% \DescribeMacro{\cauthor}
% \DescribeMacro{\eauthor}
% \begin{example}
% \cauthor{中文姓名}
% \eauthor{Your name in PinYin}
% \end{example}
%
% \myentry{申请学位名称}
% \DescribeMacro{\cdegree}
% \DescribeMacro{\edegree}
% \begin{example}
% \cdegree{您要申请什么学位}
% \edegree{degree in English}
% \end{example}
%
% \myentry{院系名称}
% \DescribeMacro{\cdepartment}
% \DescribeMacro{\edepartment}
%
% \cs{cdepartment} 可以加一个可选参数，如：\cs{cdepartmentl}\oarg{精简}\marg{详
%   细}，因为需要填写的空间有限，最好
% 给出一个详细和精简院系名称，如\textbf{计算机科学与技术}和\textbf{计算机}。
% \begin{example}
% \cdepartment[系名简称]{系名全称}
% \edepartment{Department}
% \end{example}
%
% \myentry{专业名称}
% \DescribeMacro{\cmajor}
% \DescribeMacro{\emajor}
% \begin{example}
% \cmajor{专业名称}
% \emajor{Major in English}
% \end{example}
%
% \DescribeMacro{\cfirstdiscipline}
% \DescribeMacro{\cseconddiscipline}
% \begin{example}
% \cfirstdiscipline{博士后一级学科}
% \cseconddiscipline{博士后二级学科}
% \end{example}
%
% \myentry{导师姓名}
% \DescribeMacro{\csupervisor}
% \DescribeMacro{\esupervisor}
% \begin{example}
% \csupervisor{导师~教授}
% \esupervisor{Supervisor}
% \end{example}
%
% \myentry{副导师姓名}
% \DescribeMacro{\cassosupervisor}
% \DescribeMacro{\eassosupervisor}
% \begin{example}
% \cassosupervisor{副导师~副教授}
% \eassosupervisor{Small Boss}
% \end{example}
%
% \myentry{联合导师}
% \DescribeMacro{\ccosupervisor}
% \DescribeMacro{\ecosupervisor}
% 硕士生联合指导教师，博士生联合导师。
% \begin{example}
% \ccosupervisor{联合导师~教授}
% \ecosupervisor{Tiny Boss}
% \end{example}
%
% \myentry{论文成文日期}
% \DescribeMacro{\cdate}
% \DescribeMacro{\edate}
% \DescribeMacro{\postdoctordate}
% 默认为当前时间，也可以自己指定。
% \begin{example}
% \cdate{中文日期}
% \edate{English Date}
% \postdoctordate{2009年7月——2011年7月} % 博士后研究起止日期
% \end{example}
%
% \myentry{博士后封面其它参数}
% \DescribeMacro{\catalognumber}
% \DescribeMacro{\udc}
% \DescribeMacro{\id}
% \begin{example}
% \catalognumber{分类号}
% \udc{udc}
% \id{编号}
% \end{example}
%
% \myentry{摘要}
% \DescribeEnv{cabstract}
% \DescribeEnv{eabstract}
% \begin{example}
% \begin{cabstract}
%  摘要请写在这里...
% \end{cabstract}
% \begin{eabstract}
%  here comes English abstract...
% \end{eabstract}
% \end{example}
%
% \myentry{关键词}
% \DescribeMacro{\ckeywords}
% \DescribeMacro{\ekeywords}
% 关键词用英文逗号分割写入相应的命令中，模板会解析各关键词并生成符合不同论文格式
% 要求的关键词格式。
% \begin{example}
% \ckeywords{关键词 1, 关键词 2}
% \ekeywords{keyword 1, key word 2}
% \end{example}
%
% \subsubsection{其它部分}
% \label{sec:otherparts}
% 论文其它主要部分命令：
%
% \myentry{符号对照表}
% \DescribeEnv{denotation}
% 主要符号表环境。简单定义的一个 list，跟 description 非常类似，使用方法参见示例
% 文件。带一个可选参数，用来指定符号列的宽度（默认为 2.5cm）。
% \begin{example}
% \begin{denotation}
%   \item[E] 能量
%   \item[m] 质量
%   \item[c] 光速
% \end{denotation}
% \end{example}
%
% 如果你觉得符号列的宽度不满意，那可以这样来调整：
% \begin{example}
% \begin{denotation}[1.5cm] % 设置为 1.5cm
%   \item[E] 能量
%   \item[m] 质量
%   \item[c] 光速
% \end{denotation}
% \end{example}
%
% \myentry{索引}
% 插图、表格和公式三个索引命令分别如下，将其插入到期望的位置即可（带星号的命令表
% 示对应的索引表不会出现在目录中）：
%
% \begin{center}
% \begin{tabular}{ll}
% \hline
%   {\heiti 命令} & {\heiti 说明} \\\hline
% \cs{listoffigures} & 插图索引\\
% \cs{listoffigures*} & \\\hline
% \cs{listoftables} & 表格索引\\
% \cs{listoftables*} & \\\hline
% \cs{listofequations} & 公式索引\\
% \cs{listofequations*} & \\\hline
% \end{tabular}
% \end{center}
%
% \LaTeX{} 默认支持插图和表格索引，是通过 \cs{caption} 命令完成的，因此它们必须出
% 现在浮动环境中，否则不被计数。
%
% 有的同学不想让某个表格或者图片出现在索引里面，那么请使用命令 \cs{caption*}，这
% 个命令不会给表格编号，也就是出来的只有标题文字而没有``表~xx''，``图~xx''，否则
% 索引里面序号不连续就显得不伦不类，这也是 \LaTeX{} 里星号命令默认的规则。
%
% 有这种需求的多是英文资料翻译部分，如果你觉得附录中英文原文中的表格和
% 图片显示成``表''和``图''很不协调的话，一个很好的办法还是用 \cs{caption*}，参数
% 随便自己写，具体用法请参看示例文档。
%
% 如果你的确想让它编号，但又不想让它出现在索引中的话，那就自己改一改模板的代码吧，
% 我目前不打算给模板增加这种另类命令。
%
% 公式索引为本模板扩展，模板扩展了 \pkg{amsmath} 几个内部命令，使得公式编号样式和
% 自动索引功能非常方便。一般来说，你用到的所有数学环境编号都没问题了，这个可以参
% 看示例文档。如果你有个非常特殊的数学环境需要加入公式索引，那么请使
% 用 \cs{equcaption}\marg{编号}。此命令表示 equation caption，带一个参数，即显示
% 在索引中的编号。因为公式与图表不同，我们很少给一个公式附加一个标题，之所以起这
% 么个名字是因为图表就是通过 \cs{caption} 加入索引的，\cs{equcaption} 完全就是为
% 了生成公式列表，不产生什么标题。
%
% 使用方法如下。假如有一个非 equation 数学环境 mymath，只要在其中写一
% 句 \cs{equcaption} 就可以将它加入公式列表。
% \begin{example}
% \begin{mymath}
%   \label{eq:emc2}\equcaption{\ref{eq:emc2}}
%   E=mc^2
% \end{mymath}
% \end{example}
%
% 当然 mymath 正文中公式的编号需要你自己来做。
%
% 同图表一样，附录中的公式有时候也不希望它跟全文统一编号，而且不希望它出现在公式
% 索引中，目前的解决办法就是利用 \cs{tag*}\marg{公式编号} 来解决。用法很简单，此
% 处不再罗嗦，实例请参看示例文档附录 A 的前两个公式。
%
% \myentry{简历}
% \DescribeEnv{resume}\DescribeMacro{\resumeitem}
% 开启个人简历章节，包括发表文章列表等。其实就是一个 chapter。里面的每个子项目请用命令 |\resumeitem{sub title}|。
%
% 这里就不再列举例子了，请参看示例文档的 data/resume.tex。
%
% \myentry{附录}
% \DescribeEnv{appendix}
% 所有的附录都插到这里来。因为附录会更改默认的 chapter 属性，而后面的{\heiti 个人简
%   历}又需要恢复，所以实现为环境可以保证全局的属性不受影响。
% \begin{example}
% \begin{appendix}
%  %%% Local Variables: 
%%% mode: latex
%%% TeX-master: "../main"
%%% End: 


\chapter{外文资料的调研阅读报告或书面翻译}
\section{单目标规划}
北冥有鱼，其名为鲲。鲲之大，不知其几千里也。化而为鸟，其名为鹏。鹏之背，不知其几
千里也。怒而飞，其翼若垂天之云。是鸟也，海运则将徙于南冥。南冥者，天池也。 
\begin{equation}\tag*{(123)}
 p(y|\mathbf{x}) = \frac{p(\mathbf{x},y)}{p(\mathbf{x})}=
\frac{p(\mathbf{x}|y)p(y)}{p(\mathbf{x})}
\end{equation}

吾生也有涯，而知也无涯。以有涯随无涯，殆已！已而为知者，殆而已矣！为善无近名，为
恶无近刑，缘督以为经，可以保身，可以全生，可以养亲，可以尽年。

\subsection{线性规划}
庖丁为文惠君解牛，手之所触，肩之所倚，足之所履，膝之所倚，砉然响然，奏刀騞然，莫
不中音，合于桑林之舞，乃中经首之会。
\begin{table}[ht]
\centering
  \centering
  \caption*{表~1\hskip1em 这是手动编号但不出现在索引中的一个表格例子}
  \label{tab:badtabular3}
  \begin{tabular}[c]{|c|m{0.8in}|c|c|c|c|c|}\hline
    \multicolumn{2}{|c|}{Network Topology} & \# of nodes & 
    \multicolumn{3}{c|}{\# of clients} & Server \\\hline
    GT-ITM & Waxman Transit-Stub & 600 &
    \multirow{2}{2em}{2\%}& 
    \multirow{2}{2em}{10\%}& 
    \multirow{2}{2em}{50\%}& 
    \multirow{2}{1.2in}{Max. Connectivity}\\\cline{1-3}
    \multicolumn{2}{|c|}{Inet-2.1} & 6000 & & & &\\\hline
    \multirow{2}{1in}{Xue} & Rui  & Ni &\multicolumn{4}{c|}{\multirow{2}*{\ucasthesis}}\\\cline{2-3}
    & \multicolumn{2}{c|}{ABCDEF} &\multicolumn{4}{c|}{} \\\hline
\end{tabular}  
\end{table}

文惠君曰：“嘻，善哉！技盖至此乎？”庖丁释刀对曰：“臣之所好者道也，进乎技矣。始臣之
解牛之时，所见无非全牛者；三年之后，未尝见全牛也；方今之时，臣以神遇而不以目视，
官知止而神欲行。依乎天理，批大郤，导大窾，因其固然。技经肯綮之未尝，而况大坬乎！
良庖岁更刀，割也；族庖月更刀，折也；今臣之刀十九年矣，所解数千牛矣，而刀刃若新发
于硎。彼节者有间而刀刃者无厚，以无厚入有间，恢恢乎其于游刃必有余地矣。是以十九年
而刀刃若新发于硎。虽然，每至于族，吾见其难为，怵然为戒，视为止，行为迟，动刀甚微，
謋然已解，如土委地。提刀而立，为之而四顾，为之踌躇满志，善刀而藏之。”

文惠君曰：“善哉！吾闻庖丁之言，得养生焉。”


\subsection{非线性规划}
孔子与柳下季为友，柳下季之弟名曰盗跖。盗跖从卒九千人，横行天下，侵暴诸侯。穴室枢
户，驱人牛马，取人妇女。贪得忘亲，不顾父母兄弟，不祭先祖。所过之邑，大国守城，小
国入保，万民苦之。孔子谓柳下季曰：“夫为人父者，必能诏其子；为人兄者，必能教其弟。
若父不能诏其子，兄不能教其弟，则无贵父子兄弟之亲矣。今先生，世之才士也，弟为盗
跖，为天下害，而弗能教也，丘窃为先生羞之。丘请为先生往说之。”
\begin{figure}[h]
  \centering
  \includegraphics{hello}
  \caption*{图~1\hskip1em 这是手动编号但不出现索引中的图片的例子}
  \label{tab:badfigure3}    
\end{figure}

柳下季曰：“先生言为人父者必能诏其子，为人兄者必能教其弟，若子不听父之诏，弟不受
兄之教，虽今先生之辩，将奈之何哉？且跖之为人也，心如涌泉，意如飘风，强足以距敌，
辩足以饰非。顺其心则喜，逆其心则怒，易辱人以言。先生必无往。”

孔子不听，颜回为驭，子贡为右，往见盗跖。

\subsection{整数规划}
盗跖乃方休卒徒大山之阳，脍人肝而餔之。孔子下车而前，见谒者曰：“鲁人孔丘，闻将军
高义，敬再拜谒者。”谒者入通。盗跖闻之大怒，目如明星，发上指冠，曰：“此夫鲁国之
巧伪人孔丘非邪？为我告之：尔作言造语，妄称文、武，冠枝木之冠，带死牛之胁，多辞缪
说，不耕而食，不织而衣，摇唇鼓舌，擅生是非，以迷天下之主，使天下学士不反其本，妄
作孝弟，而侥幸于封侯富贵者也。子之罪大极重，疾走归！不然，我将以子肝益昼餔之膳。”


\chapter{其它附录}
前面两个附录主要是给本科生做例子。其它附录的内容可以放到这里，当然如果你愿意，可
以把这部分也放到独立的文件中，然后将其 \verb|\input| 到主文件中。

%  \input{data/appendix02}
% \end{appendix}
% \end{example}
%
% \myentry{致谢声明}
% \DescribeEnv{ack}
% 把致谢做成一个环境更好一些，直接往里面写感谢的话就可以啦！下面是数学系一位同
% 学致谢里的话，拿过来做个广告，多希望每个人都能写这么一句啊！
% \begin{example}
% \begin{ack}
%   ……
%   还要特别感谢清华大学计算机系薛瑞尼同学在论文格式等方面给我的很多帮助！
%   顺便感谢下中科院计算所朝鲁同学对论文格式的修改！
% \end{ack}
% \end{example}
%
% \myentry{列表环境}
% \DescribeEnv{itemize}
% \DescribeEnv{enumerate}
% \DescribeEnv{description}
% 为了适合中文习惯，模板将这三个常用的列表环境用 \pkg{paralist} 对应的压缩环境替
% 换。一方面满足了多余空间的清楚，另一方面可以自己指定标签的样式和符号。细节请参
% 看 \pkg{paralist} 文档，此处不再赘述。
%
%
% \subsection{数学环境}
% \label{sec:math}
% \ucasthesis{} 定义了常用的数学环境：
%
% \begin{center}
% \begin{tabular}{*{7}{l}}\hline
%   axiom & theorem & definition & proposition & lemma & conjecture &\\
%   公理 & 定理 & 定义 & 命题 & 引理 & 猜想 &\\\hline
%   proof & corollary & example & exercise & assumption & remark & problem \\
%   证明 & 推论 & 例子& 练习 & 假设 & 注释 & 问题\\\hline
% \end{tabular}
% \end{center}
%
% 比如：
% \begin{example}
% \begin{definition}
% 道千乘之国，敬事而信，节用而爱人，使民以时。
% \end{definition}
% \end{example}
% 产生（自动编号）：\\[5pt]
% \fbox{{\heiti 定义~1.1~~~} {道千乘之国，敬事而信，节用而爱人，使民以时。}}
%
% 列举出来的数学环境毕竟是有限的，如果想用{\heiti 胡说}这样的数学环境，那么很容易定义：
% \begin{example}
% \newtheorem{nonsense}{胡说}[chapter]
% \end{example}
%
% 然后这样使用：
% \begin{example}
% \begin{nonsense}
% 契丹武士要来中原夺武林秘笈。—— 慕容博
% \end{nonsense}
% \end{example}
% 产生（自动编号）：\\[5pt]
% \fbox{{\heiti 胡说~1.1~~~} {契丹武士要来中原夺武林秘笈。\kern0.3ex\rule[0.8ex]{2em}{0.1ex}\kern0.3ex 慕容博}}
%
% \subsection{自定义以及其它}
% \label{sec:othercmd}
% 模板的配置文件 ucasthesis.cfg 中定义了很多固定词汇，一般无须修改。如果有特殊需求，
% 推荐在导言区使用 \cs{renewcommand}。当然，导言区里可以直接使用中文。
%
%
% \section{ThuThesis致谢}
% \label{sec:thanks}
% 感谢这些年来一直陪伴 ThuThesis 成长的新老同学，大家的需求是模板前
% 进的动力，大家的反馈是模板提高的机会。
% 
% 此版本加入了博士后出站报告的支持，本意为制作一个支持清华所有学位报告
% 的模板，孰料学校于近期对硕士、博士论文规范又有调整，未能及时更新，见
% 谅!
%
% 本人已于近期离开清华，虽不忍模板存此瑕疵，然精力有限，必不能如往日及
% 时升级，还望新的同学能参与或者接手，继续为大家服务。
% 
% \section{UCASThesis致谢}
% 感谢ThuThesis模板的作者，为国内学位论文模板提供了良好的实现基础。
%
% 感谢huzhenzhen编写的biblatex-gb7714-2015，方便的解决了本模板文献格式问题。
% 
% 同时基于 ThuThesis 修改的版本 UCASThesis 离不开首批试用该模板的林健博士，
% 唐帝鑫， 杨阳师兄提出的宝贵修改意见。在此，我要特别感谢林健博士为此模板
% 精心的校对、修正所付出的辛勤汗水。没有林健博士对我的督促、帮助，
% 不会有该模板的诞生。
% 希望国科大的同学们继续完善该模板，欢迎fork， 但我更希望各分支版本能够
% 发起pull request造福更多同学。 同时也希望各位同学在使用本模板时，能够对模板进行致谢。
%
% \StopEventually{\PrintChanges\PrintIndex}
% \clearpage
%
% \section{实现细节}
%
% \subsection{基本信息}
%    \begin{macrocode}
%<cls>\NeedsTeXFormat{LaTeX2e}[1999/12/01]
%<cls>\ProvidesClass{ucasthesis}
%<cfg>\ProvidesFile{ucasthesis.cfg}
%<cls|cfg>[2012/07/28 4.8dev University of Chinese Academy of Sciences Thesis Template]
%    \end{macrocode}
%
% \subsection{定义选项}
% \label{sec:defoption}
% TODO: 所有的选项用 \pkg{xkeyval} 来重构，现在的太罗唆了。
%
% 定义论文类型以及是否涉密
%    \begin{macrocode}
%<*cls>
\hyphenation{ucas-Thesis}
\def\ucasthesis{\textsc{ucasthesis}}
\def\version{4.8dev}
\newif\ifucas@master\ucas@masterfalse
\newif\ifucas@doctor\ucas@doctorfalse
\newif\ifucas@secret\ucas@secretfalse
\DeclareOption{master}{\ucas@mastertrue}
\DeclareOption{doctor}{\ucas@doctortrue}
\DeclareOption{secret}{\ucas@secrettrue}
%    \end{macrocode}
%
%
% 如果需要使用 arial 字体，请打开 [arial] 选项
%    \begin{macrocode}
\newif\ifucas@arial   %!
\DeclareOption{arial}{\ucas@arialtrue}
%    \end{macrocode}
%
% 目录中英文是否用 arial
%    \begin{macrocode}
\newif\ifucas@arialtoc
\DeclareOption{arialtoc}{\ucas@arialtoctrue}
%    \end{macrocode}
% 章节标题中的英文是否用 arial
%    \begin{macrocode}
\newif\ifucas@arialtitle
\DeclareOption{arialtitle}{\ucas@arialtitletrue}
%    \end{macrocode}
%
% noraggedbottom 选项
%    \begin{macrocode}
\newif\ifucas@raggedbottom\ucas@raggedbottomtrue
\DeclareOption{noraggedbottom}{\ucas@raggedbottomfalse}
%    \end{macrocode}
%
% 将选项传递给 ctexbook 类
%    \begin{macrocode}
\DeclareOption*{\PassOptionsToClass{\CurrentOption}{ctexbook}}
%    \end{macrocode}
%
% \cs{ExecuteOptions} 的参数之间用逗号分割，不能有空格。开始不知道，折腾了老半
% 天。
%    \begin{macrocode}
\ExecuteOptions{utf,arialtitle}%?
\ProcessOptions\relax
\RequirePackage{ifxetex}
 \ifxetex
  \LoadClass[cs4size,a4paper,openany,nocap,nofonts,UTF8]{ctexbook}
 \else
  \LoadClass[cs4size,a4paper,openany,nocap,winfonts,UTF8]{ctexbook}
\fi%    \end{macrocode}
%
% 用户至少要提供一个选项：指定论文类型。
%    \begin{macrocode}

\ifucas@master\relax\else
  \ifucas@doctor\relax\else
      \ClassError{ucasthesis}%
                   {You have to specify one of thesis options: master or doctor.}{}
  \fi
\fi

%    \end{macrocode}
%
% \subsection{装载宏包}
% \label{sec:loadpackage}
%
% 引用的宏包和相应的定义。
%    \begin{macrocode}
\RequirePackage{ifthen,calc}
%    \end{macrocode}
%
% \AmSTeX{} 宏包，用来排出更加漂亮的公式。
%    \begin{macrocode}
\RequirePackage{amsmath}
%    \end{macrocode}
%
% 用很爽的 \pkg{txfonts} 替换 \pkg{mathptmx} 宏包，同时用它自带的 typewriter 字
% 体替换 courier。必须出现在 \AmSTeX{} 之后。
%    \begin{macrocode}
\RequirePackage{txfonts}
%    \end{macrocode}
%
% 图形支持宏包。
%    \begin{macrocode}
\RequirePackage{graphicx}
%    \end{macrocode}
%
% 并排图形。\pkg{subfigure}、\pkg{subfig} 已经不再推荐，用新的 \pkg{subcaption}。
% 浮动图形和表格标题样式。\pkg{caption2} 已经不推荐使用，采用新的 \pkg{caption}。
%    \begin{macrocode}
\RequirePackage[labelformat=simple]{subcaption}
%    \end{macrocode}
%
% 为了兼容ctex 2.0+
%    \begin{macrocode}
\RequirePackage{indentfirst}
%    \end{macrocode}

% 更好的列表环境。
%    \begin{macrocode}
\RequirePackage[neverdecrease]{paralist}
\setdefaultleftmargin{3em}{2.2em}{1.87em}{1.7em}{1em}{1em}
%    \end{macrocode}
%
% raggedbottom，禁止Latex自动调整多余的页面底部空白，并保持脚注仍然在底部。
%    \begin{macrocode}
\ifucas@raggedbottom
  \RequirePackage[bottom,perpage]{footmisc}
  \raggedbottom
\else
  \RequirePackage[perpage]{footmisc}
\fi
%    \end{macrocode}
%
% 中文支持，我们使用 ctex 宏包。
%    \begin{macrocode}
\ifdefined \ctexset
\ctexset{autoindent=true}
\fi
\renewcommand{\baselinestretch}{1.0}
\RequirePackage{zhnumber}
\ifxetex
  \xeCJKsetup{AutoFakeBold=false,AutoFakeSlant=false,AllowBreakBetweenPuncts=true}
  \punctstyle{quanjiao}
  \normalspacedchars{-}
  
  \def\CJK@null{\kern\CJKnullspace\Unicode{48}{7}\kern\CJKnullspace}
  \defaultfontfeatures{Mapping=tex-text} % use TeX --
%    \end{macrocode}
% 默认采用中易的四款 (宋，黑，楷，仿宋) 免费字体。需要手工
% 修改 fontname.def 文件。缺少中文字体的 Linux 用户可以通过 fontname.def 文件定义字体。
%    \begin{macrocode}
  \input{fontname.def}

  \setmainfont{Times New Roman}
  \setsansfont{Arial}
  \setmonofont{Courier New}
\else
  \RequirePackage{CJKspace}
%    \end{macrocode}
% arial 字体需要单独安装，如果不使用 arial 字体，可以用 helvet 字体 |\textsf|
% 模拟，二者基本没有差别。
%    \begin{macrocode}
  \ifucas@arial
    \IfFileExists{arial.sty}%
                 {\RequirePackage{arial}}%
                 {\ClassWarning{ucasthesis}{no arial.sty availiable!}}
  \fi
\fi
%    \end{macrocode}
%
% 定理类环境宏包，其中 \pkg{amsmath} 选项用来兼容 \AmSTeX{} 的宏包
%    \begin{macrocode}
\RequirePackage[amsmath,thmmarks,hyperref]{ntheorem}
%    \end{macrocode}
%
% 表格控制
%    \begin{macrocode}
\RequirePackage{array}
\RequirePackage{longtable}
%    \end{macrocode}
%
% 使用三线表：\cs{toprule}，\cs{midrule}，\cs{bottomrule}。
%    \begin{macrocode}
\RequirePackage{booktabs}
%    \end{macrocode}
%
% 参考文献引用宏包。
%    \begin{macrocode}
%\RequirePackage[numbers,sort&compress]{natbib}

%    \end{macrocode}
%
% 生成有书签的 pdf 及其开关，请结合 gbk2uni 避免书签乱码。
%    \begin{macrocode}
\RequirePackage{hyperref}
\ifxetex
  \hypersetup{%
    CJKbookmarks=true}
\else
  \hypersetup{%
    unicode=true,
    CJKbookmarks=false}
\fi

\RequirePackage{refcount}
\newcommand{\chapterref}[1]{%
  \@ifundefined{hyperref}{}{\hyperref[#1]%
  {\zhnumber{\getrefnumber{#1}}}}}

\RequirePackage{xcolor}
\definecolor{ucascitecolor}{rgb}{0.0, 0.50, 0.0}
\definecolor{ucaslinkcolor}{rgb}{0.75, 0.0, 0.0}
\definecolor{ucasurlcolor}{rgb}{0.0, 0.0, 0.75}
\hypersetup{%
  bookmarksnumbered=true,
  bookmarksopen=true,
  bookmarksopenlevel=1,
  breaklinks=true,
  colorlinks=true,
  citecolor=ucascitecolor,
  linkcolor=ucaslinkcolor,
  urlcolor=ucasurlcolor,
  plainpages=false,
  pdfpagelabels,
  pdfborder=0 0 0}
%    \end{macrocode}
%
% dvips 模式下网址断字有问题，请手工加载 breakurl 这个宏包解决之。
%
% 设置 url 样式，与上下文一致
%    \begin{macrocode}
\urlstyle{same}
%</cls>
%    \end{macrocode}
%
%
% \subsection{主文档格式}
% \label{sec:mainbody}
%
% \subsubsection{Three matters}
% 我们的单面和双面模式与常规的不太一样。
%    \begin{macrocode}
%<*cls>
\renewcommand\frontmatter{%
  \scleardoublepage
  \@mainmatterfalse
  \pagenumbering{Roman}
  \pagestyle{ucas@empty}}
\renewcommand\mainmatter{%
  \fcleardoublepage
  \@mainmattertrue
  \pagenumbering{arabic}
  \pagestyle{ucas@headings}}
  %#\ifucas@bachelor\pagestyle{ucas@plain}\else\pagestyle{ucas@plain}\fi}
\renewcommand\backmatter{%
  \cleardoublepage
  %\pagenumbering{arabic}
  \pagestyle{ucas@headings}
  \@mainmattertrue}
%</cls>
%    \end{macrocode}
%
%
% \subsubsection{字体}
% \label{sec:font}
%
% 重定义字号命令
%
% Ref 1:
% \begin{verbatim}
% 参考科学出版社编写的《著译编辑手册》(1994年)
% 七号       5.25pt       1.845mm
% 六号       7.875pt      2.768mm
% 小五       9pt          3.163mm
% 五号      10.5pt        3.69mm
% 小四      12pt          4.2175mm
% 四号      13.75pt       4.83mm
% 三号      15.75pt       5.53mm
% 二号      21pt          7.38mm
% 一号      27.5pt        9.48mm
% 小初      36pt         12.65mm
% 初号      42pt         14.76mm
%
% 这里的 pt 对应的是 1/72.27 inch，也就是 TeX 中的标准 pt
% \end{verbatim}
%
% Ref 2:
% WORD 中的字号对应该关系如下:
% \begin{verbatim}
% 初号 = 42bp = 14.82mm = 42.1575pt
% 小初 = 36bp = 12.70mm = 36.135 pt
% 一号 = 26bp = 9.17mm = 26.0975pt
% 小一 = 24bp = 8.47mm = 24.09pt
% 二号 = 22bp = 7.76mm = 22.0825pt
% 小二 = 18bp = 6.35mm = 18.0675pt
% 三号 = 16bp = 5.64mm = 16.06pt
% 小三 = 15bp = 5.29mm = 15.05625pt
% 四号 = 14bp = 4.94mm = 14.0525pt
% 小四 = 12bp = 4.23mm = 12.045pt
% 五号 = 10.5bp = 3.70mm = 10.59375pt
% 小五 = 9bp = 3.18mm = 9.03375pt
% 六号 = 7.5bp = 2.56mm
% 小六 = 6.5bp = 2.29mm
% 七号 = 5.5bp = 1.94mm
% 八号 = 5bp = 1.76mm
%
% 1bp = 72.27/72 pt
% \end{verbatim}
%
% \begin{macro}{\ucas@define@fontsize}
% 根据习惯定义字号。用法：
%
% \cs{ucas@define@fontsize}\marg{字号名称}\marg{磅数}
%
% 避免了字号选择和行距的紧耦合。所有字号定义时为单倍行距，并提供选项指定行距倍数。
%    \begin{macrocode}
%<*cls>
\newlength\ucas@linespace
\newcommand{\ucas@choosefont}[2]{%
   \setlength{\ucas@linespace}{#2*\real{#1}}%
   \fontsize{#2}{\ucas@linespace}\selectfont}
\def\ucas@define@fontsize#1#2{%
  \expandafter\newcommand\csname #1\endcsname[1][\baselinestretch]{%
    \ucas@choosefont{##1}{#2}}}
%    \end{macrocode}
% \end{macro}
% \begin{macro}{\chuhao}
% \begin{macro}{\xiaochu}
% \begin{macro}{\yihao}
% \begin{macro}{\xiaoyi}
% \begin{macro}{\erhao}
% \begin{macro}{\xiaoer}
% \begin{macro}{\sanhao}
% \begin{macro}{\xiaosan}
% \begin{macro}{\sihao}
% \begin{macro}{\banxiaosi}
% \begin{macro}{\xiaosi}
% \begin{macro}{\dawu}
% \begin{macro}{\wuhao}
% \begin{macro}{\xiaowu}
% \begin{macro}{\liuhao}
% \begin{macro}{\xiaoliu}
% \begin{macro}{\qihao}
% \begin{macro}{\bahao}
%    \begin{macrocode}
\ucas@define@fontsize{chuhao}{42bp}
\ucas@define@fontsize{xiaochu}{36bp}
\ucas@define@fontsize{yihao}{26bp}
\ucas@define@fontsize{xiaoyi}{24bp}
\ucas@define@fontsize{erhao}{22bp}
\ucas@define@fontsize{xiaoer}{18bp}
\ucas@define@fontsize{sanhao}{16bp}
\ucas@define@fontsize{xiaosan}{15bp}
\ucas@define@fontsize{sihao}{14bp}
\ucas@define@fontsize{banxiaosi}{13bp}
\ucas@define@fontsize{xiaosi}{12bp}
\ucas@define@fontsize{dawu}{11bp}
\ucas@define@fontsize{wuhao}{10.5bp}
\ucas@define@fontsize{xiaowu}{9bp}
\ucas@define@fontsize{liuhao}{7.5bp}
\ucas@define@fontsize{xiaoliu}{6.5bp}
\ucas@define@fontsize{qihao}{5.5bp}
\ucas@define@fontsize{bahao}{5bp}
%    \end{macrocode}
% \end{macro}
% \end{macro}
% \end{macro}
% \end{macro}
% \end{macro}
% \end{macro}
% \end{macro}
% \end{macro}
% \end{macro}
% \end{macro}
% \end{macro}
% \end{macro}
% \end{macro}
% \end{macro}
% \end{macro}
% \end{macro}
% \end{macro}
% \end{macro}
%
% 正文小四号 (12pt) 字，行距为固定值 20 磅。
%    \begin{macrocode}
\renewcommand\normalsize{%
  \@setfontsize\normalsize{12bp}{20bp}
  \abovedisplayskip=10bp \@plus 2bp \@minus 2bp
  \abovedisplayshortskip=10bp \@plus 2bp \@minus 2bp
  \belowdisplayskip=\abovedisplayskip
  \belowdisplayshortskip=\abovedisplayshortskip}
%</cls>
%    \end{macrocode}
%
%
% \subsubsection{页面设置}
% \label{sec:layout}
% 本来这部分应该是最容易设置的，但根据格式规定出来的结果跟学校的 WORD 样例相差很
% 大，所以只能微调。
%    \begin{macrocode}
%<*cls>
\AtBeginDvi{\special{papersize=\the\paperwidth,\the\paperheight}}
\AtBeginDvi{\special{!%
      \@percentchar\@percentchar BeginPaperSize: a4
      ^^Ja4^^J\@percentchar\@percentchar EndPaperSize}}
\setlength{\textwidth}{\paperwidth}
\setlength{\textheight}{\paperheight}
\setlength\marginparwidth{0cm}
\setlength\marginparsep{0cm}
\addtolength{\textwidth}{-5cm}
\setlength{\topmargin}{1.75cm-1in}
\setlength{\oddsidemargin}{2.5cm-1in}
\setlength{\footskip}{0.6cm}
\setlength{\headsep}{0.5cm}
\addtolength{\textheight}{-5.25cm}

\setlength{\evensidemargin}{\oddsidemargin}
\setlength{\headheight}{20pt}
\setlength{\topskip}{0pt}
\setlength{\skip\footins}{15pt}
%</cls>
%    \end{macrocode}
%
% \subsubsection{页眉页脚}
% \label{sec:headerfooter}
% 新的一章从奇数页开始 (openright)，所以必须保证它前面那页如果没有内容也必须
% 没有页眉页脚。(code stolen from \pkg{fancyhdr})
%    \begin{macrocode}
%<*cls>
\let\ucas@cleardoublepage\cleardoublepage
\newcommand{\ucas@clearemptydoublepage}{%
  \clearpage{\pagestyle{ucas@sfoot}\ucas@cleardoublepage}}
\newcommand{\ucas@sclearemptydoublepage}{%
  \clearpage{\pagestyle{ucas@empty}\ucas@cleardoublepage}}
\newcommand{\ucas@fclearemptydoublepage}{%
  \clearpage{\pagestyle{ucas@foot}\ucas@cleardoublepage}}
\let\cleardoublepage\ucas@clearemptydoublepage
\let\scleardoublepage\ucas@sclearemptydoublepage
\let\fcleardoublepage\ucas@fclearemptydoublepage

%    \end{macrocode}
%
% 定义页眉和页脚。chapter 自动调用 thispagestyle{ucas@plain}，所以要重新定义 ucas@plain。
% \begin{macro}{\ps@ucas@empty}
% \begin{macro}{\ps@ucas@plain}
% \begin{macro}{\ps@ucas@headings}
% 定义三种页眉页脚格式：
% \begin{itemize}
% \item \texttt{ucas@empty}：页眉页脚都没有
% \item \texttt{ucas@plain}：只显示页脚的页码
% \item \texttt{ucas@headings}：页眉页脚同时显示
% \end{itemize}
%    \begin{macrocode}
\def\ps@ucas@empty{%
  \let\@oddhead\@empty%
  \let\@evenhead\@empty%
  \let\@oddfoot\@empty%
  \let\@evenfoot\@empty}
\def\ps@ucas@foot{%
  \let\@oddhead\@empty%
  \let\@evenhead\@empty%
  \def\@oddfoot{\hfil\xiaowu\thepage\hfil}%
  \let\@evenfoot=\@oddfoot}
\def\ps@ucas@sfoot{%
  \let\@oddhead\@empty%
  \let\@evenhead\@empty%
  \def\@oddfoot{\hfil\xiaowu\thepage}%
  \def\@evenfoot{\xiaowu\thepage\hfil}}
\def\ps@ucas@plain{%
  \def\@oddhead{\vbox to\headheight{%
    \hb@xt@\textwidth{\hfill\xiaowu\songti\leftmark\hfill}%
      \vskip2pt\hbox{\vrule width\textwidth height0.4pt depth0pt}}}
  \def\@oddfoot{\hfil\xiaowu\thepage\hfil}%
  \def\@evenhead{\vbox to\headheight{%
      \hb@xt@\textwidth{\xiaowu\songti%
       \hfill \leftmark \hfill}%
      \vskip2pt\hbox{\vrule width\textwidth height0.4pt depth0pt}}}
  \def\@oddfoot{\hfil\xiaowu\thepage\hfil}%
  \let\@evenfoot=\@oddfoot}
\def\ps@ucas@headings{%
  \def\@oddhead{\vbox to\headheight{%
    \hb@xt@\textwidth{\hfill\xiaowu\songti\leftmark\hfill}%
      \vskip2pt\hbox{\vrule width\textwidth height0.4pt depth0pt}}}
  \def\@evenhead{\vbox to\headheight{%
      \hb@xt@\textwidth{\xiaowu\songti%
       \hfill \headingtitle \hfill}%
      \vskip2pt\hbox{\vrule width\textwidth height0.4pt depth0pt}}}
  \def\@oddfoot{\hfil\xiaowu\thepage}
  \def\@evenfoot{\xiaowu\thepage\hfil}}

%    \end{macrocode}
% \end{macro}
% \end{macro}
% \end{macro}
%
% 其实可以直接写到 \cs{chapter} 的定义里面。
%    \begin{macrocode}
\renewcommand{\chaptermark}[1]{\@mkboth{\@chapapp\  ~~#1}{}}
%</cls>
%    \end{macrocode}
%
%
% \subsubsection{段落}
% \label{sec:paragraph}
%
% 段落之间的竖直距离
%    \begin{macrocode}
%<*cls>
\setlength{\parskip}{0pt \@plus2pt \@minus0pt}
%    \end{macrocode}
%
% 调整默认列表环境间的距离，以符合中文习惯。
% \begin{macro}{ucas@item@space}
%    \begin{macrocode}
\def\ucas@item@space{%
  \let\itemize\compactitem
  \let\enditemize\endcompactitem
  \let\enumerate\compactenum
  \let\endenumerate\endcompactenum
  \let\description\compactdesc
  \let\enddescription\endcompactdesc}
%</cls>
%    \end{macrocode}
% \end{macro}
%
%
% \subsubsection{脚注}
% \label{sec:footnote}
%
% 脚注字体：宋体小五，单倍行距。悬挂缩进 1.5 字符。标号在正文中是上标，在脚注中为
% 正体。默认情况下 \cs{@makefnmark} 显示为上标，同时为脚标和正文所用，所以如果要区
% 分，必须分别定义脚注的标号和正文的标号。
% \begin{macro}{\ucas@textcircled}
% 生成带圈的脚注数字。最多处理到 99，当然这个很容易扩展了。
%    \begin{macrocode}
%<*cls>
\def\ucas@textcircled#1{%
  \ifnum \value{#1} <10 \textcircled{\xiaoliu\arabic{#1}}
  \else\ifnum \value{#1} <100 \textcircled{\qihao\arabic{#1}}\fi
  \fi}
%    \end{macrocode}
% \end{macro}
%    \begin{macrocode}
\renewcommand{\thefootnote}{\ucas@textcircled{footnote}}
\renewcommand{\thempfootnote}{\ucas@textcircled{mpfootnote}}
\def\footnoterule{\vskip-3\p@\hrule\@width0.3\textwidth\@height0.4\p@\vskip2.6\p@}
\let\ucas@footnotesize\footnotesize
\renewcommand\footnotesize{\ucas@footnotesize\xiaowu[1.5]}
\def\@makefnmark{\textsuperscript{\hbox{\normalfont\@thefnmark}}}
\long\def\@makefntext#1{
  \bgroup
    \newbox\ucas@tempboxa
    \setbox\ucas@tempboxa\hbox{%
      \hb@xt@ 2em{\@thefnmark\hss}}
    \leftmargin\wd\ucas@tempboxa
    \rightmargin\z@
    \linewidth \columnwidth
    \advance \linewidth -\leftmargin
    \parshape \@ne \leftmargin \linewidth
    \footnotesize
    \@setpar{{\@@par}}%
    \leavevmode
    \llap{\box\ucas@tempboxa}%
    #1
  \par\egroup}
%</cls>
%    \end{macrocode}
%
%
% \subsubsection{数学相关}
% \label{sec:equation}
% 允许太长的公式断行、分页等。
%    \begin{macrocode}
%<*cls>
\allowdisplaybreaks[4]
\renewcommand\theequation{\ifnum \c@chapter>\z@ \thechapter-\fi\@arabic\c@equation}
%    \end{macrocode}
%
% 公式距前后文的距离由 4 个参数控制，参见 \cs{normalsize} 的定义。
%
% 公式改成 (1-1) 的形式。
% 同时为了让 \pkg{amsmath} 的 \cs{tag*} 命令得到正确的格式，我们必须修改这些代
% 码。\cs{make@df@tag} 是定义 \cs{tag*} 和 \cs{tag} 内部命令的。
% \cs{make@df@tag@@} 处理 \cs{tag*}，我们就改它！
% \begin{verbatim}
% \def\make@df@tag{\@ifstar\make@df@tag@@\make@df@tag@@@}
% \def\make@df@tag@@#1{%
%   \gdef\df@tag{\maketag@@@{#1}\def\@currentlabel{#1}}}
% \end{verbatim}
%    \begin{macrocode}
\def\make@df@tag{\@ifstar\ucas@make@df@tag@@\make@df@tag@@@}
\def\ucas@make@df@tag@@#1{\gdef\df@tag{\ucas@maketag{#1}\def\@currentlabel{#1}}}
% redefinitation of tagform brokes eqref!
\renewcommand{\eqref}[1]{\textup{(\ref{#1})}}
\renewcommand\theequation{\ifnum \c@chapter>\z@ \thechapter-\fi\@arabic\c@equation}
%\ifucas@bachelor
%  \def\ucas@maketag#1{\maketag@@@{%
%    (\ignorespaces\text{\equationname\hskip0.5em}#1\unskip\@@italiccorr)}}
%  \def\tagform@#1{\maketag@@@{%
%    (\ignorespaces\text{\equationname\hskip0.5em}#1\unskip\@@italiccorr)\equcaption{#1}}}
%\else
\def\ucas@maketag#1{\maketag@@@{(\ignorespaces #1\unskip\@@italiccorr)}}
\def\tagform@#1{\maketag@@@{(\ignorespaces #1\unskip\@@italiccorr)\equcaption{#1}}}
%\fi
%    \end{macrocode}
% ^^A 使公式编号随着每开始新的一节而重新开始。
% ^^A \@addtoreset{eqation}{section}
%
% 解决证明环境中方块乱跑的问题。
%    \begin{macrocode}
\gdef\@endtrivlist#1{%  % from \endtrivlist
  \if@inlabel \indent\fi
  \if@newlist \@noitemerr\fi
  \ifhmode
    \ifdim\lastskip >\z@ #1\unskip \par
      \else #1\unskip \par \fi
  \fi
  \if@noparlist \else
    \ifdim\lastskip >\z@
       \@tempskipa\lastskip \vskip -\lastskip
      \advance\@tempskipa\parskip \advance\@tempskipa -\@outerparskip
      \vskip\@tempskipa
    \fi
    \@endparenv
  \fi #1}
%    \end{macrocode}
%
% 定理字样使用黑体，正文使用宋体，冒号隔开
%    \begin{macrocode}
\theorembodyfont{\kaishu }
\theoremheaderfont{\heiti }
%</cls>
%<*cfg>
%%\theoremsymbol{\ensuremath{\blacksquare}}
\theoremsymbol{\ensuremath{\square}}
\theoremstyle{plain}
%\theoremstyle{nonumberplain}
\newtheorem*{proof}{证明}
\theoremsymbol{}

\newtheorem{assumption}{假设}[chapter]
\newtheorem{definition}{定义}[chapter]
\newtheorem{proposition}{命题}[chapter]
\newtheorem{lemma}{引理}[chapter]
\newtheorem{theorem}{定理}[chapter]
\newtheorem{axiom}{公理}[chapter]
\newtheorem{corollary}{推论}[chapter]
\newtheorem{exercise}{练习}[chapter]
\newtheorem{example}{例}[chapter]
\newtheorem{remark}{注释}[chapter]
\newtheorem{problem}{问题}[chapter]
\newtheorem{conjecture}{猜想}[chapter]
%</cfg>
%    \end{macrocode}
%
% \subsubsection{浮动对象以及表格}
% \label{sec:float}
% 设置浮动对象和文字之间的距离
%    \begin{macrocode}
%<*cls>
\setlength{\floatsep}{12bp \@plus4pt \@minus1pt}
\setlength{\intextsep}{12bp \@plus4pt \@minus2pt}
\setlength{\textfloatsep}{12bp \@plus4pt \@minus2pt}
\setlength{\@fptop}{0bp \@plus1.0fil}
\setlength{\@fpsep}{12bp \@plus2.0fil}
\setlength{\@fpbot}{0bp \@plus1.0fil}
%    \end{macrocode}
%
% 下面这组命令使浮动对象的缺省值稍微宽松一点，从而防止幅度对象占据过多的文本页面，
% 也可以防止在很大空白的浮动页上放置很小的图形。
%    \begin{macrocode}
\renewcommand{\textfraction}{0.15}
\renewcommand{\topfraction}{0.85}
\renewcommand{\bottomfraction}{0.65}
\renewcommand{\floatpagefraction}{0.60}
%    \end{macrocode}
%
% 定制浮动图形和表格标题样式
% \begin{itemize}
%   \item 图表标题写作五号
%   \item 去掉图表号后面的冒号。图序与图名文字之间空一个汉字符宽度。
%   \item 图：caption 在下，段前空 6 磅，段后空 12 磅
%   \item 表：caption 在上，段前空 12 磅，段后空 6 磅
% \end{itemize}
%    \begin{macrocode}
\usepackage[list=off]{bicaption}
\let\old@tabular\@tabular
\def\ucas@tabular{\wuhao[1.5]\old@tabular}
\DeclareCaptionLabelFormat{ucas}{{#1~\rmfamily #2}}
\DeclareCaptionLabelSeparator{ucas}{\hspace{1em}}
\DeclareCaptionFont{ucas}{\heiti \wuhao[1.5]}
\captionsetup{labelformat=ucas,labelsep=ucas,labelfont=ucas,font=ucas}
\captionsetup[table]{position=top,belowskip={12bp-\intextsep},aboveskip=6bp}
\captionsetup[figure]{position=bottom,belowskip={12bp-\intextsep},aboveskip=6bp}
\captionsetup[sub]{font=ucas,skip=6bp}
\captionsetup[figure][bi-second]{name=Figure} %设置图的英文编号前缀
\captionsetup[table][bi-second]{name=Table} %设置表的英文编号前缀
\renewcommand{\thesubfigure}{(\alph{subfigure})}
\renewcommand{\thesubtable}{(\alph{subtable})}
% \renewcommand{\p@subfigure}{:}
%    \end{macrocode}
% 我们采用 \pkg{longtable} 来处理跨页的表格。同样我们需要设置其默认字体为五号。
%    \begin{macrocode}
\let\ucas@LT@array\LT@array
\def\LT@array{\wuhao[1.5]\ucas@LT@array} % set default font size
%    \end{macrocode}
%
% \begin{macro}{\hlinewd}
% 简单的表格使用三线表推荐用 \cs{hlinewd}。如果表格比较复杂还是用 \pkg{booktabs} 的命
% 令好一些。
%    \begin{macrocode}
\def\hlinewd#1{%
  \noalign{\ifnum0=`}\fi\hrule \@height #1 \futurelet
    \reserved@a\@xhline}
%</cls>
%    \end{macrocode}
% \end{macro}
%
%
% \subsubsection{中文标题定义}
% \label{sec:theor}
%    \begin{macrocode}
%<*cfg>
\renewcommand\contentsname{目\hspace{1em}录}
\newcommand\scontentsname{目录}
\renewcommand\listfigurename{图目录}
\renewcommand\listtablename{表目录}
\newcommand\listequationname{公式索引}
\newcommand\equationname{公式}
\renewcommand\bibname{参考文献}
\renewcommand\indexname{索引}
\renewcommand\figurename{图}
\renewcommand\tablename{表}
\newcommand\CJKprepartname{第}
\newcommand\CJKpartname{部分}
\CTEXnumber{\ucas@thepart}{\@arabic\c@part}
\newcommand\CJKthepart{\ucas@thepart}
\newcommand\CJKprechaptername{第}
\newcommand\CJKchaptername{章}

\newcommand\CJKthechapter{\@arabic\c@chapter}
\renewcommand\chaptername{\CJKprechaptername\CJKthechapter\CJKchaptername}
\renewcommand\appendixname{附录}
\newcommand{\cabstractname}{摘\hspace{1em}要}
\newcommand{\csabstractname}{摘要}
\newcommand{\eabstractname}{Abstract}


\let\CJK@todaysave=\today
\def\CJK@todaysmall@short{\the\year 年 \the\month 月}
\def\CJK@todaysmall{\CJK@todaysmall@short \the\day 日}
\def\ucas@CJK@year{\zhdigits{\the\year}}
\def\ucas@year{\the\year}
\def\ucas@CJK@month{\zhdigits{\the\month}}
\def\ucas@month{\the\month}
\def\ucas@CJK@day{\zhdigits{\the\day}}
\def\CJK@todaybig@short{\ucas@year{}年\ucas@month{}月}
\def\CJK@todaybig{\CJK@todaybig@short{}\ucas@CJK@day{}日}
\def\CJK@today{\CJK@todaysmall}
\renewcommand\today{\CJK@today}
\newcommand\CJKtoday[1][1]{%
  \ifcase#1\def\CJK@today{\CJK@todaysave}
    \or\def\CJK@today{\CJK@todaysmall}
    \or\def\CJK@today{\CJK@todaybig}
  \fi}
%</cfg>
%    \end{macrocode}
%
%
% \subsubsection{章节标题}
% \label{sec:titleandtoc}
% 如果章节题目中的英文要使用 arial，那么就加上 \cs{sffamily}
%    \begin{macrocode}
%<*cls>
\ifucas@arialtitle
  \def\ucas@title@font{\sffamily}
\fi
%    \end{macrocode}
%
% \begin{macro}{\chapter}
% 章序号与章名之间空一个汉字符 黑体四号字，居中书写，单倍行距，段前空 24 磅，段
% 后空 18 磅。
%
%    \begin{macrocode}
\renewcommand\chapter{%
  \cleardoublepage%
  %\thispagestyle{ucas@headings}%
  \global\@topnum\z@%
  \@afterindenttrue%
  \secdef\@chapter\@schapter}
\def\@chapter[#1]#2{%
  \ifnum \c@secnumdepth >\m@ne
   \if@mainmatter
     \refstepcounter{chapter}%
     \addcontentsline{toc}{chapter}{\protect\numberline{\@chapapp}#1}%TODO: shit
   \else
     \addcontentsline{toc}{chapter}{#1}%
   \fi
  \else
    \addcontentsline{toc}{chapter}{#1}%
  \fi
  \chaptermark{#1}%
  \@makechapterhead{#2}}
\def\@makechapterhead#1{%
  \vspace*{17bp}
  {\parindent \z@ \centering
    \bfseries \csname ucas@title@font\endcsname\heiti\sihao[2.41]
    \ifnum \c@secnumdepth >\m@ne
      \@chapapp\hskip1em
    \fi
    #1\par\nobreak
    \vskip 16.5bp}}
\def\@schapter#1{%
  \@makeschapterhead{#1}
  \@afterheading}
\def\@makeschapterhead#1{%
  \vspace*{17bp}%
  {\parindent \z@ \centering
   \bfseries \csname ucas@title@font\endcsname\heiti\sihao[2.41]
     \def\@tempa{#1}
     \def\@tempb{\eabstractname}
     \ifx\@tempa\@tempb\fi

   \interlinepenalty\@M
   #1\par\nobreak
    \vskip 16.5bp}}
%    \end{macrocode}
% \end{macro}
%
% \begin{macro}{\ucas@chapter*}
%
% \cs{ucas@chapter*}\oarg{tocline}\marg{title}\oarg{header}: tocline 是出现在目录
% 中的条目，如果为空则此 chapter 不出现在目录中，如果省略表示目录出现 title；
% title 是章标题；header 是页眉出现的标题，如果忽略则取 title。通过这个宏我才真
% 正体会到 \TeX{} macro 的力量！
%    \begin{macrocode}
\newcounter{ucas@bookmark}
\def\ucas@chapter*{%
  \@ifnextchar [ % ]
    {\ucas@@chapter}
    {\ucas@@chapter@}}
\def\ucas@@chapter@#1{\ucas@@chapter[#1]{#1}}
\def\ucas@@chapter[#1]#2{%
  \@ifnextchar [ % ]
    {\ucas@@@chapter[#1]{#2}}
    {\ucas@@@chapter[#1]{#2}[]}}
\def\ucas@@@chapter[#1]#2[#3]{%
  \cleardoublepage
  \phantomsection
  \def\@tmpa{#1}
  \def\@tmpb{#3}
  \ifx\@tmpa\@empty
    \ifx\@tmpb\@empty
      \addtocounter{ucas@bookmark}\@ne
      \pdfbookmark[0]{#2}{ucaschapter.\theucas@bookmark}
    \else
      \addtocounter{ucas@bookmark}\@ne
      \pdfbookmark[0]{#3}{ucaschapter.\theucas@bookmark}
    \fi
  \else
    \addcontentsline{toc}{chapter}{#1}
  \fi
  \chapter*{#2}
  \ifx\@tmpb\@empty
    \@mkboth{#2}{#2}
  \else
    \@mkboth{#3}{#3}
  \fi}
%    \end{macrocode}
% \end{macro}
% \begin{macro}{\section}
% 一级节标题，例如：2.1  实验装置与实验方法
% 节标题序号与标题名之间空一个汉字符（下同）。
% 采用黑体小四字居左书写，行距为1.75倍行距，段前空 8 磅，段后空 8 磅。
%
%    \begin{macrocode}
\renewcommand\section{\@startsection {section}{1}{\z@}%
                     {8bp\@plus 1ex \@minus .2ex}%
                     {8bp \@plus .2ex}%
                     {\csname ucas@title@font\endcsname\heiti\xiaosi[1.75]}}
%    \end{macrocode}
% \end{macro}
%
% \begin{macro}{\subsection}
% 二级节标题，例如：2.1.1 实验装置
%    \begin{macrocode}
\renewcommand\subsection{\@startsection{subsection}{2}{\z@}%
                        {4bp\@plus 1ex \@minus .2ex}%
                        {4bp \@plus .2ex}%
                        {\csname ucas@title@font\endcsname\heiti\xiaosi[1.75]}}
%    \end{macrocode}
% \end{macro}
%
% \begin{macro}{\subsubsection}
% 三级节标题，例如： 2.1.2.1 归纳法
% 采用黑体小四号（12pt）字居左书写，行距为1.75倍，段前空 4 磅，段后空 4 磅。
%    \begin{macrocode}
\renewcommand\subsubsection{\@startsection{subsubsection}{3}{\z@}%
                           {4bp\@plus 1ex \@minus .2ex}%
                           {4bp \@plus .2ex}%
                           {\csname ucas@title@font\endcsname\heiti\xiaosi[1.75]}}
%</cls>
%    \end{macrocode}
% \end{macro}
%
%
% \subsubsection{目录格式}
% \label{sec:toc}
% 最多涉及 4 层，即: x.x.x.x。\par
% chapter(0), section(1), subsection(2), subsubsection(3)
%    \begin{macrocode}
%<*cls>

\setcounter{secnumdepth}{3}
\setcounter{tocdepth}{2}

%    \end{macrocode}
%
% 每章标题行前空 6 磅，后空 0 磅。如果使用目录项中英文要使用 Arial，那么就加上 \cs{sffamily}。
% 章节名中英文用 Arial 字体，页码仍用 Times。
% \begin{macro}{\tableofcontents}
%   目录生成命令。
%    \begin{macrocode}
\renewcommand\tableofcontents{%
  \fcleardoublepage
  \ucas@chapter*[]{\contentsname}[\contentsname]
  \xiaosi[1.5]\@starttoc{toc}\normalsize}
\ifucas@arialtoc
  \def\ucas@toc@font{\sffamily}
\fi
\def\@pnumwidth{2em} % 这个参数没用了
\def\@tocrmarg{2em}
\def\@dotsep{1} % 目录点间的距离
\def\@dottedtocline#1#2#3#4#5{%
  \ifnum #1>\c@tocdepth \else
    \vskip \z@ \@plus.2\p@
    {\leftskip #2\relax \rightskip \@tocrmarg \parfillskip -\rightskip
    \parindent #2\relax\@afterindenttrue
    \interlinepenalty\@M
    \leavevmode
    \@tempdima #3\relax
    \advance\leftskip \@tempdima \null\nobreak\hskip -\leftskip
    {\csname ucas@toc@font\endcsname #4}\nobreak
    \leaders\hbox{$\m@th\mkern \@dotsep mu\hbox{.}\mkern \@dotsep mu$}\hfill
    \nobreak{\normalfont \normalcolor #5}%
    \par\penalty\@highpenalty}%
  \fi}
\renewcommand*\l@chapter[2]{%
  \ifnum \c@tocdepth >\m@ne
    \addpenalty{-\@highpenalty}%
    \vskip 8bp \@plus\p@
    \setlength\@tempdima{4em}%
    \begingroup
      \parindent \z@ \rightskip \@pnumwidth
      \parfillskip -\@pnumwidth
      \leavevmode
      \advance\leftskip\@tempdima
      \hskip -\leftskip
      {\csname ucas@toc@font\endcsname\heiti\sihao #1} % numberline is called here, and it uses \@tempdima
      \leaders\hbox{$\m@th\mkern \@dotsep mu\hbox{.}\mkern \@dotsep mu$}\hfill
      \nobreak{\normalcolor #2}\par
      \penalty\@highpenalty
    \endgroup
  \fi}
\renewcommand*\l@section{\heiti\vspace*{1ex}\@dottedtocline{1}{1.2em}{2.1em}}
\renewcommand*\l@subsection{\vspace*{1ex}\@dottedtocline{2}{2em}{3em}}
%%\renewcommand*\l@subsubsection{\vspace*{1ex}\@dottedtocline{3}{3.5em}{3.8em}}
%</cls>
%    \end{macrocode}
% \end{macro}
%
%
% \subsubsection{封面和封底}
% \label{sec:cover}
% \begin{macro}{\ucas@define@term}
% 方便的定义封面的一些替换命令。
%    \begin{macrocode}
%<*cls>
\def\ucas@define@term#1{
  \expandafter\gdef\csname #1\endcsname##1{%
    \expandafter\gdef\csname ucas@#1\endcsname{##1}}
  \csname #1\endcsname{}}
%    \end{macrocode}
% \end{macro}
%
% \begin{macro}{\catalognumber}
% \begin{macro}{\udc}
% \begin{macro}{\id}
% \begin{macro}{\secretlevel}
% \begin{macro}{\secretyear}
% \begin{macro}{\ctitle}
% \begin{macro}{\cdegree}
% \begin{macro}{\cdepartment}
% \begin{macro}{\caffil}
% \begin{macro}{\cmajor}
% \begin{macro}{\cfirstdiscipline}
% \begin{macro}{\cseconddiscipline}
% \begin{macro}{\csubject}
% \begin{macro}{\cauthor}
% \begin{macro}{\csupervisor}
% \begin{macro}{\cassosupervisor}
% \begin{macro}{\ccosupervisor}
% \begin{macro}{\cdate}
% \begin{macro}{\postdoctordate}
% \begin{macro}{\etitle}
% \begin{macro}{\edegree}
% \begin{macro}{\edepartment}
% \begin{macro}{\eaffil}
% \begin{macro}{\emajor}
% \begin{macro}{\esubject}
% \begin{macro}{\eauthor}
% \begin{macro}{\esupervisor}
% \begin{macro}{\eassosupervisor}
% \begin{macro}{\ecosupervisor}
% \begin{macro}{\edate}
%    \begin{macrocode}
\ucas@define@term{catalognumber}
\ucas@define@term{udc}
\ucas@define@term{id}
\ucas@define@term{secretcontent}
\ucas@define@term{ctitle}
\ucas@define@term{cdegree}
\newcommand\cdepartment[2][]{\def\ucas@cdepartment@short{#1}\def\ucas@cdepartment{#2}}
\def\caffil{\cdepartment} % todo: for compatibility
\def\ucas@cdepartment@short{}
\def\ucas@cdepartment{}
\ucas@define@term{cmajor}
\def\csubject{\cmajor} % todo: for compatibility
\ucas@define@term{cfirstdiscipline}
\ucas@define@term{cseconddiscipline}
\ucas@define@term{cauthor}
\ucas@define@term{csupervisor}
\ucas@define@term{csupervisorplace}
\ucas@define@term{cassosupervisor}
\ucas@define@term{ccosupervisor}

\ucas@define@term{cdate}
\ucas@define@term{etitle}
\ucas@define@term{edegree}
\ucas@define@term{edepartment}
\def\eaffil{\edepartment} % todo: for compability


\ucas@define@term{eauthor}
\ucas@define@term{emajor}
\ucas@define@term{esupervisor}
\ucas@define@term{edepartmnent}
\ucas@define@term{edate}
%    \end{macrocode}
% \end{macro}
% \end{macro}
% \end{macro}
% \end{macro}
% \end{macro}
% \end{macro}
% \end{macro}
% \end{macro}
% \end{macro}
% \end{macro}
% \end{macro}
% \end{macro}
% \end{macro}
% \end{macro}
% \end{macro}
% \end{macro}
% \end{macro}
% \end{macro}
% \end{macro}
% \end{macro}
% \end{macro}
% \end{macro}
% \end{macro}
% \end{macro}
% \end{macro}
% \end{macro}
% \end{macro}
% \end{macro}
% \end{macro}
% \end{macro}
%
% 封面、摘要、版权、致谢格式定义。
% \begin{environment}{cabstract}
% \begin{environment}{eabstract}
% 摘要最好以环境的形式出现（否则命令的形式会导致开始结束的括号距离太远，我不喜
% 欢），这就必须让环境能够自己保存内容留待以后使用。ctt 上找到两种方法：1）使用
%  \pkg{amsmath} 中的 \cs{collect@body}，但是此宏没有定义为 long，不能直接用。
% 2）利用 \LaTeX{} 中环境和对应命令间的命名关系以及参数分隔符的特点非常巧妙地实
% 现了这个功能，其不足是不能嵌套环境。由于摘要部分经常会用到诸如 itemize 类似
% 的环境，所以我们不得不选择第一种负责的方法。以下是修改 \pkg{amsmath} 代码部分：
%    \begin{macrocode}
\long\@xp\def\@xp\collect@@body\@xp#\@xp1\@xp\end\@xp#\@xp2\@xp{%
  \collect@@body{#1}\end{#2}}
\long\@xp\def\@xp\push@begins\@xp#\@xp1\@xp\begin\@xp#\@xp2\@xp{%
  \push@begins{#1}\begin{#2}}
\long\@xp\def\@xp\addto@envbody\@xp#\@xp1\@xp{%
  \addto@envbody{#1}}
%    \end{macrocode}
%
% 使用 \cs{collect@body} 来构建摘要环境。
%    \begin{macrocode}
\newcommand{\ucas@@cabstract}[1]{\long\gdef\ucas@cabstract{#1}}
\newenvironment{cabstract}{\collect@body\ucas@@cabstract}{}
\newcommand{\ucas@@eabstract}[1]{\long\gdef\ucas@eabstract{#1}}
\newenvironment{eabstract}{\collect@body\ucas@@eabstract}{}
%    \end{macrocode}
% \end{environment}
% \end{environment}
%
% \begin{macro}{\ucas@parse@keywords}
%   不同论文格式关键词之间的分割不太相同，我们用 \cs{ckeywords} 和
%    \cs{ekeywords} 来收集关键词列表，然后用本命令来生成符合要求的格式。
%   \cs{expandafter} 都快把我整晕了。
%    \begin{macrocode}
\def\ucas@parse@keywords#1{
  \expandafter\gdef\csname ucas@#1\endcsname{} % todo: need or not?
  \expandafter\gdef\csname #1\endcsname##1{
    \@for\reserved@a:=##1\do{
      \expandafter\ifx\csname ucas@#1\endcsname\@empty\else
        \expandafter\g@addto@macro\csname ucas@#1\endcsname{\ignorespaces\csname ucas@#1@separator\endcsname}
      \fi
      \expandafter\expandafter\expandafter\g@addto@macro%
        \expandafter\csname ucas@#1\expandafter\endcsname\expandafter{\reserved@a}}}}
%    \end{macrocode}
% \end{macro}
% \begin{macro}{\ckeywords}
% \begin{macro}{\ekeywords}
% 利用 \cs{ucas@parse@keywords} 来定义，内部通过 \cs{ucas@ckeywords} 来引用。
%    \begin{macrocode}
\ucas@parse@keywords{ckeywords}
\ucas@parse@keywords{ekeywords}
%</cls>
%    \end{macrocode}
% \end{macro}
% \end{macro}
%
%    \begin{macrocode}
%<*cfg>
\def\ucas@ckeywords@separator{，}
\def\ucas@ekeywords@separator{,}
\def\ucas@catalog@number@title{分类号}
\def\ucas@id@title{编号}
\def\ucas@title@sep{：}
\def\ucas@secretcontent{秘密}
\def\ucas@schoolname{中国科学院大学}



\def\ucas@author@title{作者姓名}
\ifucas@doctor
  \def\ucas@apply{博士学位论文}
\else
  \def\ucas@apply{硕士学位论文}
\fi
\newcommand{\headingtitle}{\ucas@ctitle}

\def\ucas@department@title{研\hspace{\fill}究\hspace{\fill}所}
\def\ucas@major@title{学科专业}
\def\ucas@degree@title{学位类别}
\def\ucas@supervisor@title{指导教师}
\def\ucas@assosuper@title{副指导教师}

\def\ucas@cosuper@title{%
  \ifucas@doctor 联合导师\else \ifucas@master 联合指导教师\fi\fi}
\cdate{\CJK@todaybig@short}
\edate{\ifcase \month \or January\or February\or March\or April\or May%
       \or June\or July \or August\or September\or October\or November
       \or December\fi\unskip\ \ \the\year}
\newcommand{\ucas@authtitle}{中国科学院大学\\\vspace{1em}学位论文授权使用声明}
\newcommand{\ucas@authorization}{%
本人完全了解并同意遵守中国科学院有关保存和使用学位论文的规定，即中国科学院有权保留送交学位论文的副本，
允许该论文被查阅，可以按照学术研究公开原则和保护知识产权的原则公布该论文的全部或部分内容，
可以采用影印、缩印或其他复制手段保存、汇编本学位论文。\par
涉密及延迟公开的学位论文在解密或延迟期后适用本声明。

}
\newcommand{\ucas@authorizationaddon}{%
（保密论文在解密后适用本授权书。）}
\newcommand{\ucas@authorsig}{作者签名：}
\newcommand{\ucas@teachersig}{导师签名：}
\newcommand{\ucas@frontdate}{日期：}
\newcommand{\ucas@ckeywords@title}{关键词：}
%</cfg>
%    \end{macrocode}
%
% \subsubsection{致谢以及声明}
% \label{sec:ackanddeclare}
%
%    \begin{macrocode}
%<*cfg>
\newcommand{\ucas@ackname}{致谢}
\newcommand{\ucas@sackname}{致\hspace{1em}谢}
\newcommand{\ucas@declarename}{中国科学院大学\\\vspace{1em}研究生学位论文原创性声明}
\newcommand{\ucas@declaretext}{本人郑重声明：所呈交的学位论文是本人在导师的指导下独立进行研究工作所取得的成果。
尽我所知，除文中已经注明引用的内容外，本论文不包含任何其他个人或集体已经发表或撰写过的研究成果。
对论文所涉及的研究工作做出贡献的其他个人和集体，均已在文中以明确方式标明或致谢。}
\newcommand{\ucas@signature}{作者签名：}
\newcommand{\ucas@backdate}{日\hspace{2em}期：}
%</cfg>
%    \end{macrocode}
%
% \begin{macro}{\ucas@first@titlepage}
% 论文封面第一页！
%
% 题名使用一号黑体字，一行写不下时可分两行写，并采用 1.25 倍行距。
% 申请学位的学科门类: 小二号宋体字。
% 中文封面页边距：
%  上- 6.0 厘米，下- 5.5 厘米，左- 4.0 厘米，右- 4.0 厘米，装订线 0 厘米；
%
%    \begin{macrocode}
%<*cls>
\newcommand\ucas@underline[2][6em]{\hskip1pt\underline{\hb@xt@ #1{\hss#2\hss}}\hskip3pt}
\newlength{\ucas@title@width}
\def\ucas@put@title#1{\makebox[\ucas@title@width][s]{#1}}
\def\ucas@first@titlepage{%
  \ucas@first@titlepage@other}
\newcommand*{\getcmlength}[1]{\strip@pt\dimexpr0.035146\dimexpr#1\relax\relax}
\newcommand{\ucas@first@titlepage@other}{
  \begin{center}
    \vspace*{0cm}
    \parbox[t][1cm][t]{\textwidth}{%
      \ifucas@secret\hspace{25em}{\sihao[1]\songti 密级：\underline{\makebox[6em]{\ucas@secretcontent}}}\fi}
      \begin{figure}[H]
        \centering
        \includegraphics[width=90mm]{ucaslogo}
      \end{figure}
      \parbox[t][8cm][t]{\paperwidth-8cm}{
      \renewcommand{\baselinestretch}{1.3}
      \begin{center}    
      \yihao\heiti\textbf{\ucas@apply}
      \par\vskip 40bp
      \xiaosan\heiti\textbf{\ucas@ctitle}
      \end{center}}
%    \end{macrocode}
%
% 作者及导师信息部分使用三号仿宋字
%    \begin{macrocode}


\parbox[t][9cm][t]{\textwidth}{{\sihao\heiti
\begin{center}
\setlength{\ucas@title@width}{4em}
\setlength{\extrarowheight}{2.5ex}
\begin{tabular}{ll}

\ucas@author@title\ucas@title@sep     &\underline{\makebox[25em]{\ucas@cauthor}}\\

\ucas@supervisor@title\ucas@title@sep     &\underline{\makebox[25em]{\ucas@csupervisor}}\\

&\underline{\makebox[25em]{\ucas@csupervisorplace}}\\

\ucas@degree@title\ucas@title@sep     &\underline{\makebox[25em]{\ucas@cdegree}}\\

\ucas@major@title\ucas@title@sep     &\underline{\makebox[25em]{\ucas@cmajor}}\\

\ucas@put@title{\ucas@department@title}\ucas@title@sep     &\underline{\makebox[25em]{\ucas@cdepartment}}\\

\end{tabular}
\end{center}}}
%    \end{macrocode}
%
% 论文成文打印的日期，用阿拉伯数字
%    \begin{macrocode}
  \begin{center}
    {\vskip 3ex\sihao \heiti \ucas@cdate}
  \end{center}
\end{center}} % end of titlepage
%    \end{macrocode}
% \end{macro}
%
% \begin{macro}{\ucas@doctor@engcover}
% 研究生论文英文封面部分。
%    \begin{macrocode}
\newcommand{\ucas@engcover}{%
  \def\ucas@master@art{Master of Arts}
  \def\ucas@master@sci{Master of Science}
  \def\ucas@doctor@phi{Doctor of Philosophy}
  \newif\ifucas@professional
  \ucas@professionalfalse

  \begin{center}
    \vspace*{0.2cm}
    \parbox[t][3.2cm][t]{\paperwidth-7.2cm}{
      \renewcommand{\baselinestretch}{1.5}
      \begin{center}
        \xiaosan\sffamily\textbf{\ucas@etitle}
      \end{center}}


    \parbox[t][15cm][t]{\paperwidth-7.2cm}{
    \vspace*{2cm}
      \renewcommand{\baselinestretch}{1.4}
      \begin{center}
        \sihao
        \rmfamily \textbf{A \ifucas@master thesis \else dissertation \fi
        submitted to}\\
        \rmfamily \textbf{University of Chinese Academy of Sciences}\\
        \rmfamily \textbf{in partial fulfillment of the requirement}\\
        \rmfamily \textbf{for the degree of}\\
        \rmfamily \textbf{\ucas@edegree}\\
        \rmfamily \textbf{in \ucas@emajor}

        \sanhao\rmfamily \textbf{By}\\[3bp]
        \sihao
        \rmfamily \textbf{\ucas@eauthor}\\
        \sihao
        \rmfamily \textbf{Supervisor: Professor \ucas@esupervisor}
      \end{center}}

    \par\vspace{0.5cm}
    \parbox[t][2cm][b]{\paperwidth-7.2cm}{
    \renewcommand{\baselinestretch}{1.4}
    \begin{center}
      \sihao\rmfamily\textbf{\ucas@edepartment}\par
      \sihao\rmfamily\textbf{\ucas@edate}
    \end{center}}
  \end{center}}
%    \end{macrocode}
% \end{macro}
% \begin{macro}{\ucas@authorization@mk}
% 封面中论文授权部分。
%    \begin{macrocode}
\newcommand{\ucas@declare@mk}{%
\begin{center}\sihao\heiti{\ucas@declarename}\end{center}
  \vskip25pt\xiaosi[2.03]\par
  \ucas@declaretext\par
  \vskip1.5cm
  \xiaosi
  \hspace{13em}\ucas@signature\par
  \vskip0.4cm
  \hspace{13em}\ucas@backdate
}
\newcommand{\ucas@authorization@mk}{%
  %\vspace*{0.72cm} % shit code!
  \begin{center}\sihao\heiti{\ucas@authtitle}\end{center}
  \vskip40pt\xiaosi[2.03]\par
  \ucas@authorization\par
  %\ucas@authorizationaddon\par
  \vskip1.5cm
  \xiaosi
  \hspace{6em}\ucas@authorsig\hspace{6em}\ucas@teachersig\par
  \vskip0.4cm
  \hspace{6em}\ucas@backdate\hspace{6em}\ucas@backdate
  }
%    \end{macrocode}
% \end{macro}
%
%
% \begin{macro}{\makecover}
%    \begin{macrocode}
\newcommand{\makecover}{
  \phantomsection
  \pdfbookmark[-1]{\ucas@ctitle}{ctitle}
  \normalsize%
  \begin{titlepage}
%    \end{macrocode}
%
% 论文封面第一页！
%    \begin{macrocode}
    \ucas@first@titlepage
%    \end{macrocode}
%
%
% 研究生论文需要增加英文封面
%    \begin{macrocode}
\scleardoublepage
\clearpage\ucas@engcover

%    \end{macrocode}
%
% 授权说明
%    \begin{macrocode}
      \scleardoublepage%
      \begin{list}{}{%
        \topsep\z@%
        \listparindent\parindent%
        \parsep\parskip%
        \setlength{\itemsep}{3cm}
        \setlength{\leftmargin}{6mm}%
        \setlength{\rightmargin}{6mm}}%
      \item[]\ucas@declare@mk
      \item[]\ucas@authorization@mk%
      \end{list}%
  \end{titlepage}
%    \end{macrocode}
%
%
% 中英文摘要
%    \begin{macrocode}
  \normalsize
  \ucas@makeabstract
  \let\@tabular\ucas@tabular}
%</cls>
%    \end{macrocode}
% \end{macro}
%
% \subsubsection{摘要格式}
% \label{sec:abstractformat}
%
% \begin{macro}{\ucas@makeabstract}
% 中文摘要部分的标题为\textbf{摘要}，用黑体三号字。
%    \begin{macrocode}
%<*cls>
\newcommand{\ucas@makeabstract}{%
  \scleardoublepage
  \ucas@chapter*[]{\cabstractname}[] % no tocline
  \pagestyle{ucas@plain}
  \pagenumbering{Roman}
%    \end{macrocode}
%
% 摘要内容用小四号字书写，两端对齐，汉字用宋体，外文字用 Times New Roman 体，
% 标点符号一律用中文输入状态下的标点符号。
%    \begin{macrocode}
  \ucas@cabstract
%    \end{macrocode}
% 每个关键词之间空两个汉字符宽度， 且为悬挂缩进
%    \begin{macrocode}
  \vskip12bp
  \setbox0=\hbox{{\heiti\ucas@ckeywords@title}}
  \noindent\hangindent\wd0\hangafter1
    \box0\ucas@ckeywords
%    \end{macrocode}
%
% 英文摘要部分的标题为 \textbf{Abstract}，用 Arial 体三号字。研究生的英文摘要要求
% 非常怪异：虽然正文前的封面部分为右开，但是英文摘要要跟中文摘要连
% 续。
%    \begin{macrocode}
  \fcleardoublepage
  \ucas@chapter*[]{\eabstractname}[\eabstractname] % no tocline
%    \end{macrocode}
%
% 摘要内容用小四号 Times New Roman。
%    \begin{macrocode}
  \ucas@eabstract
%    \end{macrocode}
%
% 每个关键词之间空四个英文字符宽度
%    \begin{macrocode}
  \vskip12bp
  \setbox0=\hbox{\xiaosi\textbf{Key Words:\enskip}}
  \noindent
    \box0\ucas@ekeywords}
%</cls>
%    \end{macrocode}
% \end{macro}
%
% \subsubsection{主要符号表}
% \label{sec:denotationfmt}
% \begin{environment}{denotation}
% 主要符号对照表
%    \begin{macrocode}
%<*cfg>
\newcommand{\ucas@denotation@name}{主要符号对照表}
%</cfg>
%<*cls>
\newenvironment{denotation}[1][2.5cm]{
  \ucas@chapter*[]{\ucas@denotation@name} % no tocline
  \noindent\begin{list}{}%
    {\vskip-30bp\xiaosi[1.6]
     \renewcommand\makelabel[1]{##1\hfil}
     \setlength{\labelwidth}{#1} % 标签盒子宽度
     \setlength{\labelsep}{0.5cm} % 标签与列表文本距离
     \setlength{\itemindent}{0cm} % 标签缩进量
     \setlength{\leftmargin}{\labelwidth+\labelsep} % 左边界
     \setlength{\rightmargin}{0cm}
     \setlength{\parsep}{0cm} % 段落间距
     \setlength{\itemsep}{0cm} % 标签间距
    \setlength{\listparindent}{0cm} % 段落缩进量
    \setlength{\topsep}{0pt} % 标签与上文的间距
   }}{\end{list}}
%</cls>
%    \end{macrocode}
% \end{environment}
%
%

% \begin{environment}{ack}
%    \begin{macrocode}
%<*cls>

\newenvironment{ack}{%
    \cleardoublepage
    \ucas@chapter*[\ucas@ackname]{\ucas@sackname}[\ucas@sackname]
  }

%</cls>
%    \end{macrocode}
% \end{environment}
%
% \subsubsection{索引部分}
% \label{sec:threeindex}
% \begin{macro}{\listoffigures}
% \begin{macro}{\listoffigures*}
% \begin{macro}{\listoftables}
% \begin{macro}{\listoftables*}
%    \begin{macrocode}
%<*cls>
\def\ucas@starttoc#1{% #1: float type, prepend type name in \listof*** entry.
  \let\oldnumberline\numberline
  \def\numberline##1{\oldnumberline{\csname #1name\endcsname \hskip.4em ##1}}
  \@starttoc{\csname ext@#1\endcsname}
  \let\numberline\oldnumberline}
\def\ucas@listof#1{% #1: float type
  \fcleardoublepage
  \@ifstar
    {\ucas@chapter*[]{\csname list#1name\endcsname}\ucas@starttoc{#1}}
    {\ucas@chapter*[]{\csname list#1name\endcsname}\ucas@starttoc{#1}}}
\renewcommand\listoffigures{\ucas@listof{figure}}
\renewcommand*\l@figure{\vspace*{1ex}\@dottedtocline{1}{0em}{4em}}
\renewcommand\listoftables{\ucas@listof{table}}
\let\l@table\l@figure
%    \end{macrocode}
% \end{macro}
% \end{macro}
% \end{macro}
% \end{macro}
%
% \begin{macro}{\equcaption}
%   本命令只是为了生成公式列表，所以这个 caption 是假的。如果要编号最好用
%    equation 环境，如果是其它编号环境，请手动添加添加 \cs{equcaption}。
% 用法如下：
%
% \cs{equcaption}\marg{counter}
%
% \marg{counter} 指定出现在索引中的编号，一般取 \cs{theequation}，如果你是用
%  \pkg{amsmath} 的 \cs{tag}，那么默认是 \cs{tag} 的参数；除此之外可能需要你
% 手工指定。
%
%    \begin{macrocode}
\def\ext@equation{loe}
\def\equcaption#1{%
  \addcontentsline{\ext@equation}{equation}%
                  {\protect\numberline{#1}}}
%    \end{macrocode}
% \end{macro}
%
% \begin{macro}{\listofequations}
% \begin{macro}{\listofequations*}
% \LaTeX{}默认没有公式索引，此处定义自己的 \cs{listofequations}。
%    \begin{macrocode}
\newcommand\listofequations{\ucas@listof{equation}}
\let\l@equation\l@figure
%</cls>
%    \end{macrocode}
% \end{macro}
% \end{macro}
%
%
% \subsubsection{参考文献}
% \label{sec:ref}
%
% \begin{macro}{\onlinecite}
% 正文引用模式。依赖于 \pkg{natbib} 宏包，修改其中的命令。
%    \begin{macrocode}
%<*cls>
%\bibpunct{[}{]}{,}{s}{}{,}
%\renewcommand\NAT@citesuper[3]{\ifNAT@swa%
%  \unskip\kern\p@\textsuperscript{\NAT@@open #1\NAT@@close}%
%  \if*#3*\else\ (#3)\fi\else #1\fi\endgroup}
%\DeclareRobustCommand\onlinecite{\@onlinecite}
%\def\@onlinecite#1{\begingroup\let\@cite\NAT@citenum\citep{#1}\endgroup}
%\renewcommand{\cite}[1]{{\onlinecite{#1}}}
%    \end{macrocode}
% \end{macro}
%
% 参考文献的正文部分用五号字。
% 行距采用固定值 16 磅，段前空 3 磅，段后空 0 磅。
%
% \begin{macro}{\ucasdot}
% 研究生参考文献条目最后可加点，图书文献一般不加。
%    \begin{macrocode}
\def\ucasdot{\unskip.}
%    \end{macrocode}
% \end{macro}
% \begin{macro}{ucasmasterbib}
% \begin{macro}{ucasphdbib}
%   研究生模板要求外文硕士论文参考文献显示``[Master Thesis]''，而博士模板
%   则于 2007 年冬要求显示为``[M]''。对应的外文博士论文参考文献分别显示为``[Phd
%   Thesis]''和``[D]''。
%   研究生写作指南(201109)要求：
%   中文硕士学位论文标注``[硕士学位论文]''，
%   中文博士学位论文标注``[博士学位论文]''，外文学位论文标注``[D]''。
%    \begin{macrocode}
\def\ucasmasterbib{[硕士学位论文]}
\def\ucasphdbib{[博士学位论文]}
%    \end{macrocode}
% \end{macro}
% \end{macro}
% \begin{environment}{thebibliography}
% 修改默认的 thebibliography 环境，增加一些调整代码。
%    \begin{macrocode}
\renewenvironment{thebibliography}[1]{%
   \ucas@chapter*{\bibname}%
   \erhao[1.5]
   \list{\@biblabel{\@arabic\c@enumiv}}%
        {\renewcommand{\makelabel}[1]{##1\hfill}
         \settowidth\labelwidth{1.1cm}
         \setlength{\labelsep}{0.4em}
         \setlength{\itemindent}{0pt}
         \setlength{\leftmargin}{\labelwidth+\labelsep}
         \addtolength{\itemsep}{-0.7em}
         \usecounter{enumiv}%
         \let\p@enumiv\@empty
         \renewcommand\theenumiv{\@arabic\c@enumiv}}%
    \sloppy\frenchspacing
    \clubpenalty4000
    \@clubpenalty \clubpenalty
    \widowpenalty4000%
    \interlinepenalty4000%
    \sfcode`\.\@m}
   {\def\@noitemerr
     {\@latex@warning{Empty `thebibliography' environment}}%
    \endlist\frenchspacing}
%</cls>
%    \end{macrocode}
% \end{environment}
%
%
% \subsubsection{附录}
% \label{sec:appendix}
%
% \begin{environment}{appendix}
%    \begin{macrocode}
%<*cls>

\let\ucas@appendix\appendix
\renewenvironment{appendix}{%
  \ucas@appendix
  \gdef\@chapapp{\appendixname}
  %\renewcommand\theequation{\ifnum \c@chapter>\z@ \thechapter-\fi\@arabic\c@equation}
  }{}
%</cls>
%    \end{macrocode}
% \end{environment}
%
% \subsubsection{个人简历}
%
% 定义个人简历章节标题
% \begin{environment}{resume}
% 个人简历发表文章等。
%    \begin{macrocode}
%<*cls>
\newenvironment{resume}[1][\ucas@resume@title]{%
  \cleardoublepage
  \ucas@chapter*{#1}}{}
%</cls>
%    \end{macrocode}
% \end{environment}
%
% \begin{macro}{\resumeitem}
% 个人简历里面会出现的以发表文章，在投文章等。
%    \begin{macrocode}
%<*cfg>
\newcommand{\ucas@resume@title}{作者简历及攻读学位期间发表的学术论文与研究成果}

%</cfg>
%<*cls>
\newcommand{\resumeitem}[1]{\vspace{2.5em}{\noindent\xiaosi\heiti {#1}： \par}}
%</cls>
%    \end{macrocode}
% \end{macro}
%
% \subsubsection{书脊}
% \label{sec:spine}
% \begin{macro}{\spine}
% 单独使用书脊命令会在新的一页产生竖排书脊。
%    \begin{macrocode}
%<*cls>
\newcommand{\spine}[1][\ucas@ctitle]{
  \newpage\thispagestyle{empty}\heiti\xiaosi
  \parbox[t][23cm][t]{1em}{#1\vfill\ucas@cauthor\vfill 中国科学院大学}}
%</cls>
%    \end{macrocode}
% \end{macro}
%
% \subsubsection{索引}
%
% 生成索引的一些命令，虽然我们暂时还用不到。
%    \begin{macrocode}
%<*cls>
\iffalse
\newcommand{\bs}{\symbol{'134}}%Print backslash
% \newcommand{\bs}{\ensuremath{\mathtt{\backslash}}}%Print backslash
% Index entry for a command (\cih for hidden command index
\newcommand{\cih}[1]{%
  \index{commands!#1@\texttt{\bs#1}}%
  \index{#1@\texttt{\hspace*{-1.2ex}\bs #1}}}
\newcommand{\ci}[1]{\cih{#1}\texttt{\bs#1}}
% Package
\newcommand{\pai}[1]{%
  \index{packages!#1@\textsf{#1}}%
  \index{#1@\textsf{#1}}%
  \textsf{#1}}
% Index entry for an environment
\newcommand{\ei}[1]{%
  \index{environments!\texttt{#1}}%
  \index{#1@\texttt{#1}}%
  \texttt{#1}}
% Indexentry for a word (Word inserted into the text)
\newcommand{\wi}[1]{\index{#1}#1}
\fi
%</cls>
%    \end{macrocode}
%
% \subsubsection{自定义命令和环境}
% \label{sec:userdefine}
%
%
% \subsubsection{其它}
% \label{sec:other}
%
% 在模板文档结束时即装入配置文件，这样用户就能在导言区进行相应的修改，否则
% 必须在 document 开始后才能，感觉不好。
%    \begin{macrocode}
%<*cls>
\AtEndOfClass{%%
%% This is file `ucasthesis.cfg',
%% generated with the docstrip utility.
%%
%% The original source files were:
%%
%% ucasthesis.dtx  (with options: `cfg')
%% 
%% This is a generated file.
%% 
%% Copyright (C) 2005-2018 by Xue Ruini <xueruini@gmail.com>
%% Copyright (C) 2013-2018 derived by Chaolu <wuchaolu@gmail.com>
%% 
%% This file may be distributed and/or modified under the
%% conditions of the LaTeX Project Public License, either version 1.3a
%% of this license or (at your option) any later version.
%% The latest version of this license is in:
%% 
%% http://www.latex-project.org/lppl.txt
%% 
%% and version 1.3a or later is part of all distributions of LaTeX
%% version 2004/10/01 or later.
%% 
%% 
%% This is the configuration file of the ucasthesis package with LaTeX2e.
%% 

\ProvidesFile{ucasthesis.cfg}
[2012/07/28 4.8dev University of Chinese Academy of Sciences Thesis Template]
%%\theoremsymbol{\ensuremath{\blacksquare}}
\theoremsymbol{\ensuremath{\square}}
\theoremstyle{plain}
\newtheorem*{proof}{证明}
\theoremsymbol{}

\newtheorem{assumption}{假设}[chapter]
\newtheorem{definition}{定义}[chapter]
\newtheorem{proposition}{命题}[chapter]
\newtheorem{lemma}{引理}[chapter]
\newtheorem{theorem}{定理}[chapter]
\newtheorem{axiom}{公理}[chapter]
\newtheorem{corollary}{推论}[chapter]
\newtheorem{exercise}{练习}[chapter]
\newtheorem{example}{例}[chapter]
\newtheorem{remark}{注释}[chapter]
\newtheorem{problem}{问题}[chapter]
\newtheorem{conjecture}{猜想}[chapter]
\renewcommand\contentsname{目\hspace{1em}录}
\newcommand\scontentsname{目录}
\renewcommand\listfigurename{图目录}
\renewcommand\listtablename{表目录}
\newcommand\listequationname{公式索引}
\newcommand\equationname{公式}
\renewcommand\bibname{参考文献}
\renewcommand\indexname{索引}
\renewcommand\figurename{图}
\renewcommand\tablename{表}
\newcommand\CJKprepartname{第}
\newcommand\CJKpartname{部分}
\CTEXnumber{\ucas@thepart}{\@arabic\c@part}
\newcommand\CJKthepart{\ucas@thepart}
\newcommand\CJKprechaptername{第}
\newcommand\CJKchaptername{章}

\newcommand\CJKthechapter{\@arabic\c@chapter}
\renewcommand\chaptername{\CJKprechaptername\CJKthechapter\CJKchaptername}
\renewcommand\appendixname{附录}
\newcommand{\cabstractname}{摘\hspace{1em}要}
\newcommand{\csabstractname}{摘要}
\newcommand{\eabstractname}{Abstract}

\let\CJK@todaysave=\today
\def\CJK@todaysmall@short{\the\year 年 \the\month 月}
\def\CJK@todaysmall{\CJK@todaysmall@short \the\day 日}
\def\ucas@CJK@year{\zhdigits{\the\year}}
\def\ucas@year{\the\year}
\def\ucas@CJK@month{\zhdigits{\the\month}}
\def\ucas@month{\the\month}
\def\ucas@CJK@day{\zhdigits{\the\day}}
\def\CJK@todaybig@short{\ucas@year{}年\ucas@month{}月}
\def\CJK@todaybig{\CJK@todaybig@short{}\ucas@CJK@day{}日}
\def\CJK@today{\CJK@todaysmall}
\renewcommand\today{\CJK@today}
\newcommand\CJKtoday[1][1]{%
  \ifcase#1\def\CJK@today{\CJK@todaysave}
    \or\def\CJK@today{\CJK@todaysmall}
    \or\def\CJK@today{\CJK@todaybig}
  \fi}
\def\ucas@ckeywords@separator{，}
\def\ucas@ekeywords@separator{,}
\def\ucas@catalog@number@title{分类号}
\def\ucas@id@title{编号}
\def\ucas@title@sep{：}
\def\ucas@secretcontent{秘密}
\def\ucas@schoolname{中国科学院大学}

\def\ucas@author@title{作者姓名}
\ifucas@doctor
  \def\ucas@apply{博士学位论文}
\else
  \def\ucas@apply{硕士学位论文}
\fi
\newcommand{\headingtitle}{\ucas@ctitle}

\def\ucas@department@title{研\hspace{\fill}究\hspace{\fill}所}
\def\ucas@major@title{学科专业}
\def\ucas@degree@title{学位类别}
\def\ucas@supervisor@title{指导教师}
\def\ucas@assosuper@title{副指导教师}

\def\ucas@cosuper@title{%
  \ifucas@doctor 联合导师\else \ifucas@master 联合指导教师\fi\fi}
\cdate{\CJK@todaybig@short}
\edate{\ifcase \month \or January\or February\or March\or April\or May%
       \or June\or July \or August\or September\or October\or November
       \or December\fi\unskip\ \ \the\year}
\newcommand{\ucas@authtitle}{中国科学院大学\\\vspace{1em}学位论文授权使用声明}
\newcommand{\ucas@authorization}{%
本人完全了解并同意遵守中国科学院有关保存和使用学位论文的规定，即中国科学院有权保留送交学位论文的副本，
允许该论文被查阅，可以按照学术研究公开原则和保护知识产权的原则公布该论文的全部或部分内容，
可以采用影印、缩印或其他复制手段保存、汇编本学位论文。\par
涉密及延迟公开的学位论文在解密或延迟期后适用本声明。

}
\newcommand{\ucas@authorizationaddon}{%
（保密论文在解密后适用本授权书。）}
\newcommand{\ucas@authorsig}{作者签名：}
\newcommand{\ucas@teachersig}{导师签名：}
\newcommand{\ucas@frontdate}{日期：}
\newcommand{\ucas@ckeywords@title}{关键词：}
\newcommand{\ucas@ackname}{致谢}
\newcommand{\ucas@sackname}{致\hspace{1em}谢}
\newcommand{\ucas@declarename}{中国科学院大学\\\vspace{1em}研究生学位论文原创性声明}
\newcommand{\ucas@declaretext}{本人郑重声明：所呈交的学位论文是本人在导师的指导下独立进行研究工作所取得的成果。
尽我所知，除文中已经注明引用的内容外，本论文不包含任何其他个人或集体已经发表或撰写过的研究成果。
对论文所涉及的研究工作做出贡献的其他个人和集体，均已在文中以明确方式标明或致谢。}
\newcommand{\ucas@signature}{作者签名：}
\newcommand{\ucas@backdate}{日\hspace{2em}期：}
\newcommand{\ucas@denotation@name}{主要符号对照表}

\newcommand{\ucas@resume@title}{作者简历及攻读学位期间发表的学术论文与研究成果}

\endinput
%%
%% End of file `ucasthesis.cfg'.
}
%    \end{macrocode}
%
% \begin{macro}{\ucas@setup@pdfinfo}
% 设置一些 pdf 文档信息，依赖于 \pkg{hyperref} 宏包。
%    \begin{macrocode}
\def\ucas@setup@pdfinfo{%
  \hypersetup{%
    pdftitle={\ucas@ctitle},
    pdfauthor={\ucas@cauthor},
    pdfsubject={\ucas@cdegree},
    pdfkeywords={\ucas@ckeywords},
    pdfcreator={\ucas@cauthor},
    pdfproducer={\ucasthesis}}}
%    \end{macrocode}
% \end{macro}
%
% 应用对列表环境的修改。
%    \begin{macrocode}
\AtEndOfClass{\sloppy\ucas@item@space}
%</cls>
%    \end{macrocode}
%
% \Finale
%
% \iffalse
%    \begin{macrocode}
%<*dtx-style>
\ProvidesPackage{dtx-style}

\RequirePackage{calc}
\RequirePackage{array,longtable}
\RequirePackage{fancybox,fancyvrb}
\RequirePackage{xcolor}
\RequirePackage{ifxetex}

\ifxetex
  \RequirePackage[nofonts,UTF8,hyperref]{ctex}
  \input{fontname.def}
\else
  \RequirePackage[winfonts,UTF8,hyperref]{ctex}
  \RequirePackage{txfonts}
\fi
\RequirePackage{hyperref}
\ifxetex
  \hypersetup{%
    CJKbookmarks=true}
\else
  \hypersetup{%
    unicode=true,
    CJKbookmarks=false}
\fi
\hypersetup{%
  bookmarksnumbered=true,
  bookmarksopen=true,
  bookmarksopenlevel=1,
  breaklinks=true,
  linkcolor=blue,
  plainpages=false,
  pdfpagelabels,
  pdfborder=0 0 0}
\RequirePackage{url}
\RequirePackage{indentfirst}

\setlength{\parskip}{4pt plus1pt minus0pt}
\setlength{\topsep}{0pt}
\setlength{\partopsep}{0pt}
\setlength{\parindent}{20pt}
\addtolength{\oddsidemargin}{-1cm}
\advance\textwidth 1.5cm
\addtolength{\topmargin}{-1cm}
\addtolength{\headsep}{0.3cm}
\addtolength{\textheight}{2.3cm}

\renewcommand{\baselinestretch}{1.3}
\setlength{\shadowsize}{3pt}
\def\DescribeOption#1{\SpecialOptionIndex{#1}}
\def\SpecialOptionIndex#1{\index{#1\actualchar\textbf{#1}}}
\renewenvironment{description}
  {\list{}{\setlength\labelwidth{2cm}%
           \setlength\labelsep{3pt}%
           \setlength\leftmargin{\labelwidth+\labelsep}%
           \addtolength{\itemsep}{3pt}%
           \renewcommand\makelabel[1]{%
             \shadowbox{\color{blue!90}\texttt##1}\DescribeOption{##1}}}
  }{\endlist}
\DefineVerbatimEnvironment{example}{Verbatim}%
  {frame=single,framerule=0.3mm,rulecolor=\color{red!75!green!50!blue},%
   fillcolor=\color{red!75!green!50!blue!15},framesep=2mm,baselinestretch=1.2,%
   fontsize=\small,gobble=1}
\DefineVerbatimEnvironment{shell}{Verbatim}%
  {frame=single,framerule=0.3mm,rulecolor=\color{red!85!green!60},%
   fillcolor=\color{red!85!green!10},framesep=2mm,fontsize=\small,gobble=1}
\long\def\myentry#1{\vskip5pt\par\noindent\llap{{\color{blue}\fangsong #1}}\marginpar{\strut}\hskip\parindent}
\def\tableofcontents{\renewcommand{\baselinestretch}{1.0}\@starttoc{toc}}
\def\DescribeMacro{\Describe@Macro}
\def\Describe@Macro#1{\PrintDescribeMacro{#1}\SpecialUsageIndex{#1}}
\def\PrintDescribeMacro#1{{\color{-red!75!green!50!blue!55}\MacroFont \string #1\hskip1em}}
\def\ps@headings{%
  \let\@oddfoot\@empty
  \def\@oddhead{\vbox{%
    \hb@xt@ \textwidth{\llap{\fbox{\rightmark\rule[-2pt]{0pt}{13pt}}}\hfil\thepage}%
    \vskip-0.7pt%
    \hb@xt@ \textwidth{\hrulefill}}}
  \let\@evenfoot\@oddfoot
  \let\@evenhead\@oddhead
  \let\@mkboth\markboth
  \def\sectionmark##1{%
    \markright{\ifnum \c@secnumdepth >\m@ne
      \thesection\quad
      \fi
      ##1}}
  \def\subsectionmark##1{%
    \markright{\ifnum \c@secnumdepth >\m@ne
      \thesubsection\quad
      \fi
      ##1}}
  \def\subsubsectionmark##1{%
    \markright{\ifnum \c@secnumdepth >\m@ne
      \thesubsubsection\quad
      \fi
      ##1}}}
\renewcommand\section{\@startsection{section}{1}{\z@}%
                                    {-3.5ex \@plus -1ex \@minus -.2ex}%
                                    {2.3ex \@plus.2ex}%
                                    {\normalfont\Large\bfseries}}

\renewcommand\subsection{\@startsection{subsection}{2}{\z@}%
                                       {-3.25ex\@plus -1ex \@minus -.2ex}%
                                       {1.5ex \@plus .2ex}%
                                       {\normalfont\large\bfseries}}
\renewcommand\subsubsection{\@startsection{subsubsection}{3}{\z@}%
                                          {-3.25ex\@plus -1ex \@minus -.2ex}%
                                          {1.5ex \@plus .2ex}%
                                          {\normalfont\normalsize\bfseries}}
\renewcommand\paragraph{\@startsection{paragraph}{4}{\z@}%
                                      {3.25ex \@plus1ex \@minus.2ex}%
                                      {-1em}%
                                      {\normalfont\normalsize\bfseries}}
\renewcommand\subparagraph{\@startsection{subparagraph}{5}{\parindent}%
                                         {3.25ex \@plus1ex \@minus .2ex}%
                                         {-1em}%
                                         {\normalfont\normalsize\bfseries}}
\pagestyle{empty}
%</dtx-style>
%    \end{macrocode}
% \fi
%
\endinput
